\documentclass[11pt,a4paper,oneside]{report}

\usepackage[T1]{fontenc}
\usepackage[utf8]{inputenc}
\usepackage{lmodern}
\usepackage[margin=24mm,headheight=15pt]{geometry}
\usepackage{amsmath,amssymb,amsthm,mathtools,bm}
\usepackage{booktabs,longtable,tabularx,array,multirow}
\usepackage{graphicx}
\usepackage{xcolor}
\usepackage{microtype}
\usepackage{enumitem}
\usepackage{setspace}
\usepackage{caption,subcaption,float,placeins}
\usepackage{pdflscape}
\usepackage{fancyhdr}
\usepackage{titlesec}
\usepackage[most]{tcolorbox}
\usepackage{listings}
\usepackage{xurl}
\usepackage[colorlinks=true,linkcolor=blue!48!black,citecolor=blue!48!black,urlcolor=blue!48!black,
  pdftitle={The Geometric Snowball of Speciation: Full Derivations and Interpretations},
  pdfauthor={Daniel Ortiz-Barrientos},
  pdfsubject={Mathematical derivation, geometric interpretation, biological interpretation, and computational audit of the geometric snowball of speciation},
  pdfkeywords={speciation, reproductive isolation, quantitative genetics, G matrix, fitness curvature, principal angles, Dobzhansky-Muller incompatibilities}]{hyperref}
\usepackage{bookmark}
\usepackage[nameinlink,noabbrev]{cleveref}

\setstretch{1.08}
\setlength{\parindent}{0pt}
\setlength{\parskip}{0.58em}
\setlength{\emergencystretch}{2em}
\setcounter{tocdepth}{2}
\setcounter{secnumdepth}{3}
\numberwithin{equation}{chapter}

\definecolor{ink}{HTML}{252525}
\definecolor{deepblue}{HTML}{205493}
\definecolor{teal}{HTML}{237C79}
\definecolor{forest}{HTML}{2E6F51}
\definecolor{amber}{HTML}{A96816}
\definecolor{brick}{HTML}{9A3C31}
\definecolor{violet}{HTML}{6A4C93}
\definecolor{paleBlue}{HTML}{EDF4FB}
\definecolor{paleTeal}{HTML}{EDF8F7}
\definecolor{paleGreen}{HTML}{EFF7F1}
\definecolor{paleAmber}{HTML}{FCF6EA}
\definecolor{paleRed}{HTML}{FAEFED}
\definecolor{paleViolet}{HTML}{F4F0F8}

\pagestyle{fancy}
\fancyhf{}
\fancyhead[L]{\small The geometric snowball of speciation}
\fancyhead[R]{\small Full derivations and interpretations}
\fancyfoot[C]{\thepage}
\renewcommand{\headrulewidth}{0.4pt}

\titleformat{\chapter}[display]
  {\normalfont\bfseries\color{deepblue}}
  {\filleft\Large\chaptername\ \thechapter}
  {0.4ex}
  {\titlerule\vspace{1ex}\Huge}
\titlespacing*{\chapter}{0pt}{-16pt}{24pt}
\titleformat{\section}{\Large\bfseries\color{teal}}{\thesection}{0.7em}{}
\titleformat{\subsection}{\large\bfseries\color{forest}}{\thesubsection}{0.7em}{}
\titleformat{\subsubsection}{\normalsize\bfseries\color{amber}}{\thesubsubsection}{0.7em}{}

\newtheorem{theorem}{Theorem}[chapter]
\newtheorem{proposition}[theorem]{Proposition}
\newtheorem{corollary}[theorem]{Corollary}
\newtheorem{lemma}[theorem]{Lemma}
\newtheorem{definition}[theorem]{Definition}
\newtheorem{assumption}[theorem]{Assumption}
\newtheorem{remark}[theorem]{Remark}

\tcbset{enhanced,breakable,boxrule=0.7pt,arc=1.3mm,left=1.8mm,right=1.8mm,top=1.4mm,bottom=1.4mm}
\newtcolorbox{mathteach}[1][]{colback=paleTeal,colframe=teal,title={Mathematical teaching},fonttitle=\bfseries,#1}
\newtcolorbox{geombox}[1][]{colback=paleBlue,colframe=deepblue,title={Geometric interpretation},fonttitle=\bfseries,#1}
\newtcolorbox{biobox}[1][]{colback=paleGreen,colframe=forest,title={Biological interpretation},fonttitle=\bfseries,#1}
\newtcolorbox{connectionbox}[1][]{colback=paleAmber,colframe=amber,title={Connection to earlier results},fonttitle=\bfseries,#1}
\newtcolorbox{cautionbox}[1][]{colback=paleRed,colframe=brick,title={Assumption, limitation, or failure mode},fonttitle=\bfseries,#1}
\newtcolorbox{summarybox}[1][]{colback=paleViolet,colframe=violet,title={Result to retain},fonttitle=\bfseries,#1}
\newtcolorbox{derivationbox}[1][]{colback=white,colframe=deepblue,title={Line-by-line derivation},fonttitle=\bfseries,#1}

\newcommand{\G}{\mathbf{G}}
\newcommand{\Hh}{\mathbf{H}}
\newcommand{\Sop}{\mathbf{S}}
\newcommand{\Rr}{\mathbf{R}}
\newcommand{\Pp}{\mathbf{P}}
\newcommand{\Vv}{\mathbf{V}}
\newcommand{\Ww}{\mathbf{W}}
\newcommand{\Qq}{\mathbf{Q}}
\newcommand{\M}{\mathbf{M}}
\newcommand{\Aa}{\mathbf{A}}
\newcommand{\Bb}{\mathbf{B}}
\newcommand{\Ii}{\mathbf{I}}
\newcommand{\Zero}{\mathbf{0}}
\newcommand{\zbar}{\bar{\mathbf z}}
\newcommand{\thetav}{\boldsymbol{\theta}}
\newcommand{\betav}{\boldsymbol{\beta}}
\newcommand{\psiv}{\boldsymbol{\psi}}
\newcommand{\muv}{\boldsymbol{\mu}}
\newcommand{\cK}{\mathcal K}
\newcommand{\cT}{\mathcal T}
\newcommand{\cA}{\mathcal A}
\newcommand{\cD}{\mathcal D}
\newcommand{\R}{\mathbb R}
\newcommand{\E}{\mathbb E}
\newcommand{\Var}{\operatorname{Var}}
\newcommand{\Cov}{\operatorname{Cov}}
\newcommand{\tr}{\operatorname{tr}}
\newcommand{\rank}{\operatorname{rank}}
\newcommand{\diag}{\operatorname{diag}}
\newcommand{\svec}{\operatorname{svec}}
\newcommand{\mat}{\operatorname{mat}}
\newcommand{\Span}{\operatorname{span}}
\newcommand{\Grass}{\operatorname{Gr}}
\newcommand{\sym}{\operatorname{sym}}
\newcommand{\normF}[1]{\left\lVert #1\right\rVert_F}
\newcommand{\normtwo}[1]{\left\lVert #1\right\rVert_2}
\newcommand{\innerF}[2]{\left\langle #1,#2\right\rangle_F}
\newcommand{\dd}{\mathrm d}
\newcommand{\half}{\tfrac12}

\newcommand{\inwords}{\par\noindent\textbf{In words.}\ }
\newcommand{\why}{\par\noindent\textbf{Why this step is valid.}\ }
\newcommand{\usual}{\par\noindent\textbf{Usual linear-algebra interpretation.}\ }
\newcommand{\bio}{\par\noindent\textbf{Biological reading.}\ }
\newcommand{\linkback}{\par\noindent\textbf{Link to the previous result.}\ }

\lstset{basicstyle=\ttfamily\small,frame=single,breaklines=true,columns=fullflexible,showstringspaces=false}

\title{\textbf{The Geometric Snowball of Speciation}\\[0.6em]
\Large Every derivation step, the mathematical machinery behind it, and its biological and geometric meaning}
\author{Daniel Ortiz-Barrientos\\\small Mathematical development and computational audit}
\date{5 August 2026}

\begin{document}
\hypersetup{pageanchor=false}
\maketitle

\begin{abstract}
This document develops the geometric snowball of speciation from first principles. It begins with the geometry of multivariate quantitative genetics, derives the symmetric response operator
\[
\Sop=\G^{1/2}\Hh\G^{1/2},
\]
shows how its leading invariant subspaces define additive genetic channels, and then derives how those subspaces rotate when populations experience different deterministic or stochastic evolutionary forces. The central local result is
\[
d_C^2(t)=t^2\normF{\cK}^2+O(t^3),
\]
where \(d_C\) is the chordal distance between channels and \(\cK\) is the curvature- and gap-weighted cross-channel forcing in a fixed shared representation. A common-metric generalization is then derived for evolving \(\G_i\), including both rotation of the \(\Sop_i\) eigenspaces and deformation by the changing square-root maps \(\G_i^{1/2}\). The document then proves that channel angle is not, by itself, reproductive isolation. It becomes a hybrid-fitness load only when genetic covariance and local fitness curvature are coadapted. Under a channel-ridge model, the covariance component of hybrid log-load is exactly
\[
L_{\mathrm{cov}}=\frac{\Delta g\,\Delta h}{4}\sum_{j=1}^k\sin^2\psi_j.
\]
Finally, the theory is connected to standing genetic variation, the evolution of \(\G\), effective gene flow, and the later combinatorial accumulation of Dobzhansky--Muller incompatibilities. Every substantial algebraic step is written out and immediately translated into ordinary mathematical language, geometric meaning, biological meaning, and its relation to earlier results. The final chapters connect the framework to covariance-tensor analyses and the existing \textit{Senecio lautus} programme.
\end{abstract}
\clearpage
\hypersetup{pageanchor=true}
\pagenumbering{roman}

% Chapters are kept in separate source files for review and maintenance.
\chapter*{How to use this document}
\addcontentsline{toc}{chapter}{How to use this document}

The argument has five logically distinct layers. Keeping them separate prevents the most common conceptual errors.

\begin{enumerate}[leftmargin=1.6em]
\item \textbf{Local adaptation.} A quadratic fitness surface and the additive covariance matrix generate the response operator \(\Sop\).
\item \textbf{Channel geometry.} Eigenvectors and invariant subspaces of \(\Sop\) define the locally efficient directions of multivariate response.
\item \textbf{Channel divergence.} Perturbation theory determines how quickly two populations' channel subspaces rotate apart.
\item \textbf{Fitness calibration.} A separate hybrid-fitness model determines whether and how channel divergence creates reproductive isolation.
\item \textbf{Evolutionary handoff.} Reduced gene flow can allow lineage-specific substitutions and intrinsic incompatibilities to accumulate later.
\end{enumerate}

Each important derivation is followed by four kinds of explanation. A \emph{mathematical teaching} box explains the familiar mathematical machinery. A \emph{geometric interpretation} box turns the algebra into a picture. A \emph{biological interpretation} box states what the result means for adaptation or speciation. A \emph{connection} box shows how the result follows from, extends, or constrains something derived earlier.

\begin{cautionbox}[title={What ``all derivation steps'' means here}]
The document derives all steps that are specific to this theory, including the matrix identities on which the main biological claims depend. For foundational results whose complete proof would require a separate numerical-linear-algebra text---most notably the global Davis--Kahan theorem---the exact theorem is stated, its assumptions are unpacked, and the local formula actually used by the theory is derived in full.
\end{cautionbox}

\begin{table}[H]
\centering
\caption*{Status of the main claims}
\begin{tabularx}{\textwidth}{@{}p{0.28\textwidth}X@{}}
\toprule
Claim & Mathematical status \\
\midrule
Quadratic local fitness model & A second-order local approximation near a fitness maximum. \\
Whitened response operator \(\Sop\) & Exact when \(\G\) is treated as fixed over the local relaxation interval; also defines the instantaneous generalized response modes when \(\G\) evolves. \\
Principal-angle and projector identities & Exact. \\
Local quadratic channel snowball & Exact asymptotic result under smooth deterministic divergence and a nonzero subspace gap; the common-metric coefficient includes the direct effect of evolving \(\G^{1/2}\). \\
Linear drift baseline & Exact leading expectation under Brownian motion in local normal coordinates; the quadratic term of the observed curve also contains second-order stochastic and manifold corrections. \\
Gaussian hybrid-fitness formula & Exact under Gaussian phenotypes and quadratic fitness. \\
Covariance--curvature trace identity & Exact given the stated hybrid covariance decomposition. \\
Channel-ridge fitness theorem & Exact under the structured \(\G_i\) and \(\Hh_i\) model. \\
Geometric-to-DMI handoff & A transparent dynamical closure, not an identity or universal population-genetic law. \\
\bottomrule
\end{tabularx}
\end{table}

\tableofcontents
\clearpage
\pagenumbering{arabic}

\part{The mathematical language of the theory}

\chapter{Trait space, vectors, matrices, and coordinate systems}

\section{A phenotype is a point in a multivariate space}

Let an organism be described by \(n\) quantitative traits. We collect those trait values into a column vector
\begin{equation}
\mathbf z=
\begin{bmatrix}
z_1\\ z_2\\ \vdots\\ z_n
\end{bmatrix}
\in\R^n.
\end{equation}

\inwords The phenotype is represented by one point in an \(n\)-dimensional trait space. Each coordinate records one measured trait. A population mean \(\zbar\) is another point in the same space.

A displacement from a reference phenotype \(\thetav\) is
\begin{equation}
\mathbf d=\mathbf z-\thetav.
\end{equation}

\inwords Subtraction converts two locations into a direction and distance. The vector \(\mathbf d\) points from the reference phenotype to the observed phenotype.

\begin{geombox}
For two traits, trait space is an ordinary plane. For three traits, it is three-dimensional space. In many traits, the drawing is no longer literal, but all the same operations remain valid: points, vectors, lengths, angles, ellipsoids, and projections.
\end{geombox}

\begin{biobox}
The reference point will usually be a local phenotypic optimum. The displacement \(\mathbf d\) therefore records both how far the population is from its optimum and the combination of traits in which it is displaced.
\end{biobox}

\section{Matrices are transformations, not just tables of numbers}

An \(n\times n\) matrix \(\mathbf A\) acts on a vector by matrix multiplication:
\begin{equation}
\mathbf y=\mathbf A\mathbf x.
\end{equation}

\inwords The matrix takes an input direction \(\mathbf x\), stretches or compresses it, possibly rotates or shears it, and returns a new vector \(\mathbf y\).

The product of two matrices means composition:
\begin{equation}
\mathbf C=\mathbf A\mathbf B,
\qquad
\mathbf C\mathbf x=\mathbf A(\mathbf B\mathbf x).
\end{equation}

\inwords The rightmost transformation acts first. The product \(\mathbf A\mathbf B\) means ``apply \(\mathbf B\), then apply \(\mathbf A\).''

\begin{mathteach}
This order matters. Even if \(\mathbf A\) and \(\mathbf B\) are individually symmetric, generally \(\mathbf A\mathbf B\neq\mathbf B\mathbf A\). The matrices commute only when applying them in either order has the same effect. For symmetric matrices, commutation is equivalent to the existence of a common orthonormal eigenbasis.
\end{mathteach}

\section{Changing trait units}

Suppose we replace the original traits by a new linear coordinate system
\begin{equation}
\mathbf z'=\mathbf T\mathbf z,
\end{equation}
where \(\mathbf T\) is nonsingular. This can represent rescaling, changing units, rotating to principal components, or taking any invertible linear combination of traits.

A covariance matrix transforms as
\begin{equation}
\G'=\mathbf T\G\mathbf T^\top.
\label{eq:Gtransform}
\end{equation}

\begin{derivationbox}
Let \(\mathbf z\) have covariance \(\G\). Then
\begin{align}
\G'
&=\Cov(\mathbf z')\\
&=\Cov(\mathbf T\mathbf z)\\
&=\E\left[(\mathbf T\mathbf z-\mathbf T\E\mathbf z)
(\mathbf T\mathbf z-\mathbf T\E\mathbf z)^\top\right]\\
&=\E\left[\mathbf T(\mathbf z-\E\mathbf z)(\mathbf z-\E\mathbf z)^\top\mathbf T^\top\right]\\
&=\mathbf T\E\left[(\mathbf z-\E\mathbf z)(\mathbf z-\E\mathbf z)^\top\right]\mathbf T^\top\\
&=\mathbf T\G\mathbf T^\top.
\end{align}
Every equality simply moves a fixed linear transformation outside an expectation.
\end{derivationbox}

A quadratic-form matrix transforms differently. If fitness loss is the same biological scalar before and after the coordinate change,
\begin{equation}
\mathbf d^\top\Hh\mathbf d
=(\mathbf d')^\top\Hh'\mathbf d',
\end{equation}
with \(\mathbf d'=\mathbf T\mathbf d\). Substituting \(\mathbf d=\mathbf T^{-1}\mathbf d'\) gives
\begin{align}
\mathbf d^\top\Hh\mathbf d
&=(\mathbf T^{-1}\mathbf d')^\top\Hh(\mathbf T^{-1}\mathbf d')\\
&=(\mathbf d')^\top\mathbf T^{-\top}\Hh\mathbf T^{-1}\mathbf d'.
\end{align}
Therefore
\begin{equation}
\Hh'=\mathbf T^{-\top}\Hh\mathbf T^{-1}.
\label{eq:Htransform}
\end{equation}

\begin{geombox}
Covariance and curvature transform in complementary ways. Covariance follows the coordinates forward; curvature follows them backward so that the scalar quadratic loss remains unchanged. This complementarity is why the product \(\G\Hh\) transforms by similarity:
\[
\G'\Hh'=\mathbf T\G\Hh\mathbf T^{-1}.
\]
Similarity preserves eigenvalues and maps phenotype-space eigendirections by \(\mathbf u' = \mathbf T\mathbf u\).
\end{geombox}

\begin{cautionbox}
A Euclidean angle between raw eigenvectors is not generally invariant to arbitrary rescaling of traits. Later we construct a common-metric comparison that is invariant under \emph{every} nonsingular linear change of coordinates.
\end{cautionbox}

\chapter{Quadratic forms, ellipsoids, and the meanings of \(\G\) and \(\Hh\)}

\section{The quadratic form}

For a symmetric matrix \(\mathbf A\), the scalar
\begin{equation}
q_{\mathbf A}(\mathbf x)=\mathbf x^\top\mathbf A\mathbf x
\end{equation}
is called a quadratic form.

Writing the multiplication explicitly,
\begin{equation}
\mathbf x^\top\mathbf A\mathbf x
=\sum_{i=1}^n\sum_{j=1}^n x_iA_{ij}x_j.
\end{equation}

\inwords Diagonal terms \(A_{ii}x_i^2\) weight movement along individual trait axes. Off-diagonal terms \(A_{ij}x_ix_j\) describe how pairs of traits contribute jointly.

If \(\mathbf A\succ0\), then \(\mathbf x^\top\mathbf A\mathbf x>0\) for every nonzero \(\mathbf x\). Such a matrix is symmetric positive definite.

\section{Eigenvectors turn a quadratic form into independent axes}

Because a real symmetric matrix has an orthonormal eigendecomposition,
\begin{equation}
\mathbf A=\mathbf Q\boldsymbol\Lambda\mathbf Q^\top,
\qquad
\boldsymbol\Lambda=\diag(\lambda_1,\ldots,\lambda_n),
\qquad
\mathbf Q^\top\mathbf Q=\Ii,
\end{equation}
we may write \(\mathbf x=\mathbf Q\mathbf y\). Then
\begin{align}
\mathbf x^\top\mathbf A\mathbf x
&=(\mathbf Q\mathbf y)^\top
(\mathbf Q\boldsymbol\Lambda\mathbf Q^\top)
(\mathbf Q\mathbf y)\\
&=\mathbf y^\top\mathbf Q^\top\mathbf Q
\boldsymbol\Lambda
\mathbf Q^\top\mathbf Q\mathbf y\\
&=\mathbf y^\top\boldsymbol\Lambda\mathbf y\\
&=\sum_{j=1}^n\lambda_jy_j^2.
\label{eq:quadDiagonal}
\end{align}

\inwords Rotating into the eigenvector coordinates removes all cross-products. The quadratic form becomes a weighted sum of independent squared coordinates.

\begin{mathteach}
An eigenvector is a direction that a transformation does not turn. It is only stretched or compressed. Its eigenvalue gives the amount and sign of that scaling. Symmetry guarantees that the full collection of eigenvectors can be chosen orthonormally.
\end{mathteach}

\section{The additive covariance matrix \(\G\)}

The additive genetic covariance matrix is
\begin{equation}
\G=\Cov(\mathbf a),
\end{equation}
where \(\mathbf a\) is the vector of additive breeding values.

For any unit trait combination \(\mathbf u\), the additive variance in that direction is
\begin{equation}
V_A(\mathbf u)=\Var(\mathbf u^\top\mathbf a)
=\mathbf u^\top\G\mathbf u.
\label{eq:directionalVariance}
\end{equation}

\begin{derivationbox}
Because \(\mathbf u^\top\mathbf a\) is a scalar,
\begin{align}
\Var(\mathbf u^\top\mathbf a)
&=\mathbf u^\top\Cov(\mathbf a)\mathbf u\\
&=\mathbf u^\top\G\mathbf u.
\end{align}
The first equality is the standard rule for the variance of a linear combination; the second substitutes the definition of \(\G\).
\end{derivationbox}

The level set
\begin{equation}
\mathbf d^\top\G^{-1}\mathbf d=c
\end{equation}
is an ellipsoid. Its principal axes are the eigenvectors of \(\G\), and the squared semi-axis lengths are proportional to the eigenvalues of \(\G\).

\begin{geombox}
The \(\G\)-ellipsoid is a map of available additive variation. Long directions contain abundant additive variance; short directions contain little. The leading eigenvector is \(\mathbf g_{\max}\), the line of greatest additive variance.
\end{geombox}

\section{The curvature matrix \(\Hh\)}

At a local maximum of log fitness, the Hessian of log fitness is negative definite. We therefore define
\begin{equation}
\Hh=-\nabla^2\log W(\thetav),
\end{equation}
so that \(\Hh\succ0\). The local fitness-loss approximation is
\begin{equation}
-\left[\log W(\thetav+\mathbf d)-\log W(\thetav)\right]
\approx\half\mathbf d^\top\Hh\mathbf d.
\end{equation}

For a unit direction \(\mathbf u\), the directional curvature is
\begin{equation}
h(\mathbf u)=\mathbf u^\top\Hh\mathbf u.
\end{equation}

\inwords Large \(h(\mathbf u)\) means fitness falls sharply when phenotype moves in direction \(\mathbf u\). Small \(h(\mathbf u)\) means a shallow ridge or weak penalty.

The level sets
\begin{equation}
\mathbf d^\top\Hh\mathbf d=c
\end{equation}
are equal-fitness-loss ellipsoids. Their longest geometric axis corresponds to the smallest eigenvalue of \(\Hh\), because weak curvature allows a larger displacement at the same loss.

\begin{biobox}
The matrices \(\G\) and \(\Hh\) answer different questions. \(\G\) asks, ``Where can the population vary and respond?'' \(\Hh\) asks, ``Where does a displacement matter most for fitness?'' The channel theory begins when those two answers are combined.
\end{biobox}

\begin{figure}[H]
\centering
\includegraphics[width=\textwidth]{figures/pedagogy_G_H_S.pdf}
\caption{Three geometries used throughout the theory. The genetic covariance matrix \(\G\) defines a variance ellipsoid, the curvature matrix \(\Hh\) defines equal-fitness-loss contours, and whitening combines them into the symmetric operator \(\Sop\).}
\label{fig:GHSgeometry}
\end{figure}

\chapter{Square roots, whitening, trace, determinant, and singular values}

\section{The symmetric square root of an SPD matrix}

If
\begin{equation}
\G=\mathbf Q\diag(g_1,\ldots,g_n)\mathbf Q^\top,
\qquad g_j>0,
\end{equation}
then the symmetric square root is
\begin{equation}
\G^{1/2}=\mathbf Q\diag(\sqrt{g_1},\ldots,\sqrt{g_n})\mathbf Q^\top.
\end{equation}

Multiplying it by itself gives
\begin{align}
\G^{1/2}\G^{1/2}
&=\mathbf Q\diag(\sqrt{g_j})\mathbf Q^\top
\mathbf Q\diag(\sqrt{g_j})\mathbf Q^\top\\
&=\mathbf Q\diag(g_j)\mathbf Q^\top\\
&=\G.
\end{align}

\inwords The square root applies half of the stretching represented by \(\G\). Applying that half-transformation twice reconstructs the full covariance transformation.

The inverse square root is
\begin{equation}
\G^{-1/2}=\mathbf Q\diag(g_1^{-1/2},\ldots,g_n^{-1/2})\mathbf Q^\top.
\end{equation}

\section{Whitening}

Define a whitened displacement
\begin{equation}
\mathbf x=\G^{-1/2}\mathbf d.
\label{eq:whiteningDef}
\end{equation}
Equivalently,
\begin{equation}
\mathbf d=\G^{1/2}\mathbf x.
\label{eq:unwhiteningDef}
\end{equation}

If \(\mathbf d\) has covariance \(\G\), then
\begin{align}
\Cov(\mathbf x)
&=\Cov(\G^{-1/2}\mathbf d)\\
&=\G^{-1/2}\G\G^{-1/2}\\
&=\Ii.
\end{align}

\inwords Whitening measures every displacement in units of the available genetic standard deviation. After whitening, the genetic covariance ellipsoid becomes a sphere.

Substituting \cref{eq:unwhiteningDef} into the fitness quadratic gives
\begin{align}
\mathbf d^\top\Hh\mathbf d
&=(\G^{1/2}\mathbf x)^\top\Hh(\G^{1/2}\mathbf x)\\
&=\mathbf x^\top\G^{1/2}\Hh\G^{1/2}\mathbf x.
\end{align}
This motivates the operator
\begin{equation}
\Sop=\G^{1/2}\Hh\G^{1/2}.
\label{eq:SdefinitionEarly}
\end{equation}

\begin{geombox}
A unit vector in whitened space represents one genetic standard deviation of displacement. Therefore \(\mathbf x^\top\Sop\mathbf x\) is the fitness curvature experienced by a genetically standardised displacement. The operator \(\Sop\) is the fitness landscape viewed through the amount and correlation structure of standing genetic variation.
\end{geombox}

\section{Trace as total curvature-weighted variance}

The trace is the sum of diagonal elements:
\begin{equation}
\tr(\mathbf A)=\sum_iA_{ii}.
\end{equation}
It is invariant under cyclic permutation:
\begin{equation}
\tr(\mathbf A\mathbf B)=\tr(\mathbf B\mathbf A).
\end{equation}

For a zero-mean random vector \(\boldsymbol\varepsilon\) with covariance \(\Pp\),
\begin{align}
\E[\boldsymbol\varepsilon^\top\Hh\boldsymbol\varepsilon]
&=\E\left[\tr(\boldsymbol\varepsilon^\top\Hh\boldsymbol\varepsilon)\right]\\
&=\E\left[\tr(\Hh\boldsymbol\varepsilon\boldsymbol\varepsilon^\top)\right]\\
&=\tr\left(\Hh\E[\boldsymbol\varepsilon\boldsymbol\varepsilon^\top]\right)\\
&=\tr(\Hh\Pp).
\label{eq:traceQuadraticExpectation}
\end{align}

\inwords The trace adds the variance in every direction after weighting each direction by fitness curvature. This identity is the mathematical bridge from hybrid covariance to hybrid fitness load.

\section{Determinant as volume scaling}

For a square transformation \(\mathbf A\), \(|\det\mathbf A|\) is the factor by which \(\mathbf A\) changes volume. For an SPD covariance matrix, \(\det\G\) is the generalized variance: it measures the squared volume scale of the covariance ellipsoid.

This volume interpretation later explains why
\begin{equation}
\det(\Vv_1^\top\Vv_2)^2
\end{equation}
measures squared overlap of two channel subspaces.

\section{Singular values}

Any rectangular matrix has a singular value decomposition
\begin{equation}
\mathbf A=\mathbf U\boldsymbol\Sigma\mathbf V^\top,
\end{equation}
where \(\mathbf U\) and \(\mathbf V\) have orthonormal columns and \(\boldsymbol\Sigma\) contains nonnegative singular values.

\usual Eigenanalysis asks for directions preserved by one square operator. Singular-value analysis asks how strongly one orthonormal set projects onto another. This is why principal angles between subspaces are defined using an SVD rather than an eigenanalysis.

\begin{summarybox}
The square root \(\G^{1/2}\) converts genetically standardised coordinates back to phenotype space. The inverse square root whitens genetic covariance. The trace converts a covariance into an expected quadratic loss. The determinant measures volume. Singular values measure mutual projection between subspaces. These four pieces of machinery carry nearly the entire theory.
\end{summarybox}

\part{Deriving additive genetic channels}

\chapter{From a local fitness surface to multivariate selection}

\section{Second-order Taylor expansion at a local optimum}

Let \(f(\mathbf z)=\log W(\mathbf z)\). A second-order Taylor expansion around \(\thetav\) is
\begin{equation}
f(\thetav+\mathbf d)
=f(\thetav)
+\nabla f(\thetav)^\top\mathbf d
+\half\mathbf d^\top\nabla^2f(\thetav)\mathbf d
+R_3(\mathbf d),
\label{eq:TaylorLogFitness}
\end{equation}
where \(R_3(\mathbf d)=O(\normtwo{\mathbf d}^3)\) if the required derivatives are continuous.

At a local optimum,
\begin{equation}
\nabla f(\thetav)=\Zero.
\end{equation}
At a strict local maximum, \(\nabla^2f(\thetav)\) is negative definite. Define
\begin{equation}
\Hh=-\nabla^2f(\thetav)\succ0.
\end{equation}
Then \cref{eq:TaylorLogFitness} becomes
\begin{equation}
\log W(\thetav+\mathbf d)
=\log W_0-\half\mathbf d^\top\Hh\mathbf d+R_3(\mathbf d),
\label{eq:quadraticLogFitnessFull}
\end{equation}
where \(W_0=W(\thetav)\).

\inwords The constant term is peak log fitness. The linear term disappears because the gradient is zero at the optimum. The first nonzero local change is therefore quadratic. The minus sign ensures that moving away from the maximum lowers fitness.

\begin{cautionbox}
The quadratic landscape is local. It does not assume that the entire adaptive landscape is globally Gaussian or that a single optimum describes all phenotypes. It says that sufficiently near any smooth local maximum, the leading fitness loss is quadratic.
\end{cautionbox}

\section{Differentiating the quadratic form}

We now derive the selection gradient. Write
\begin{equation}
q(\mathbf d)=\half\mathbf d^\top\Hh\mathbf d.
\end{equation}
In coordinates,
\begin{equation}
q(\mathbf d)=\half\sum_{i=1}^n\sum_{j=1}^n d_iH_{ij}d_j.
\end{equation}
Differentiate with respect to coordinate \(d_k\):
\begin{align}
\frac{\partial q}{\partial d_k}
&=\half\sum_{j=1}^n H_{kj}d_j
+\half\sum_{i=1}^n d_iH_{ik}.
\end{align}
Because \(\Hh\) is symmetric, \(H_{ik}=H_{ki}\), so the two sums are equal:
\begin{align}
\frac{\partial q}{\partial d_k}
&=\sum_{j=1}^nH_{kj}d_j\\
&=(\Hh\mathbf d)_k.
\end{align}
Collecting all coordinates gives
\begin{equation}
\nabla_{\mathbf d}\left(\half\mathbf d^\top\Hh\mathbf d\right)
=\Hh\mathbf d.
\label{eq:quadraticGradientIdentity}
\end{equation}

Since \(\mathbf d=\zbar-\thetav\) and the optimum is locally fixed,
\begin{equation}
\frac{\partial\mathbf d}{\partial\zbar}=\Ii.
\end{equation}
Therefore the gradient of log fitness with respect to the population mean is
\begin{align}
\betav
&=\nabla_{\zbar}\log W(\zbar)\\
&=-\nabla_{\mathbf d}\left(\half\mathbf d^\top\Hh\mathbf d\right)\\
&=-\Hh\mathbf d\\
&=\Hh(\thetav-\zbar).
\label{eq:selectionGradient}
\end{align}

\inwords Selection points back toward the optimum. Curvature changes both the magnitude and the direction of that pull. If the fitness surface is tilted relative to the measured trait axes, selection on one displacement component can generate a gradient involving several traits.

\begin{geombox}
The vector \(-\Hh\mathbf d\) is normal to the local equal-fitness contour passing through \(\mathbf d\). It points in the direction of steepest increase in log fitness under the ordinary Euclidean metric.
\end{geombox}

\section{Lande's response equation closes the loop}

The continuous-time multivariate response equation is
\begin{equation}
\dot{\zbar}=\G\betav.
\label{eq:LandeContinuous}
\end{equation}
Substituting \cref{eq:selectionGradient},
\begin{align}
\dot{\zbar}
&=\G\Hh(\thetav-\zbar)\\
&=-\G\Hh(\zbar-\thetav).
\end{align}
Define the mean displacement
\begin{equation}
\mathbf d=\zbar-\thetav.
\end{equation}
If \(\thetav\) is fixed over the local interval,
\begin{equation}
\dot{\mathbf d}=\dot{\zbar}.
\end{equation}
Hence
\begin{equation}
\boxed{\dot{\mathbf d}=-\G\Hh\mathbf d.}
\label{eq:localDynamicsGH}
\end{equation}

\begin{derivationbox}
The logic is sequential.
\begin{enumerate}[leftmargin=1.6em]
\item The displacement \(\mathbf d\) determines the local selection gradient through \(-\Hh\mathbf d\).
\item The additive covariance matrix transforms that gradient into an evolutionary response through \(\G\betav\).
\item Their composition is \(-\G\Hh\mathbf d\).
\end{enumerate}
No new biological assumption is introduced in the matrix product beyond the quadratic local landscape and the standard linear response equation.
\end{derivationbox}

\begin{connectionbox}
When \(\Hh=h\Ii\), selection has the same curvature in every direction. Then
\[
\dot{\mathbf d}=-h\G\mathbf d.
\]
The response directions are the eigenvectors of \(\G\), recovering the classical lines-of-least-resistance setting. The geometric snowball begins by relaxing the isotropic-curvature restriction.
\end{connectionbox}

\section{Moving optima}

If the optimum changes through time, then
\begin{equation}
\mathbf d(t)=\zbar(t)-\thetav(t)
\end{equation}
and
\begin{equation}
\dot{\mathbf d}=\dot{\zbar}-\dot{\thetav}.
\end{equation}
Therefore
\begin{equation}
\dot{\mathbf d}=-\G\Hh\mathbf d-\dot{\thetav}.
\label{eq:movingOptimumDynamics}
\end{equation}

\inwords The first term is the population's restoring response. The second term is environmental forcing caused by movement of the optimum. A population tracks well when its response rates are fast relative to the speed and direction of optimum movement.

\begin{biobox}
Different habitats can diverge through different optima, different curvature matrices, different covariance matrices, or all three. Mean divergence is controlled strongly by the optima and forcing. Channel divergence concerns the geometry of the response operator and can occur even when current means are similar.
\end{biobox}

\chapter{Whitening the dynamics and deriving the symmetric channel operator}

\section{Why \(\G\Hh\) is mathematically awkward}

The local dynamics in \cref{eq:localDynamicsGH} involve \(\G\Hh\). Although \(\G\) and \(\Hh\) are symmetric,
\begin{equation}
(\G\Hh)^\top=\Hh\G,
\end{equation}
which generally differs from \(\G\Hh\). Thus \(\G\Hh\) is usually nonsymmetric.

\inwords A nonsymmetric operator can have nonorthogonal right eigenvectors, distinct left eigenvectors, and a geometry that depends awkwardly on the coordinate metric. The biological dynamics are valid, but the raw representation is not the cleanest one for measuring rotations and angles.

\section{Whitened dynamics when \(\G\) is locally fixed}

Define
\begin{equation}
\mathbf x=\G^{-1/2}\mathbf d.
\end{equation}
If \(\G\) is constant over the local relaxation interval,
\begin{align}
\dot{\mathbf x}
&=\G^{-1/2}\dot{\mathbf d}\\
&=\G^{-1/2}(-\G\Hh\mathbf d)\\
&=-\G^{1/2}\Hh\mathbf d\\
&=-\G^{1/2}\Hh\G^{1/2}\mathbf x.
\end{align}
The third line uses
\begin{equation}
\G^{-1/2}\G=\G^{1/2},
\end{equation}
which follows immediately from the eigendecomposition of \(\G\). The fourth line substitutes \(\mathbf d=\G^{1/2}\mathbf x\).

Define
\begin{equation}
\boxed{\Sop=\G^{1/2}\Hh\G^{1/2}.}
\label{eq:channelOperator}
\end{equation}
Then
\begin{equation}
\boxed{\dot{\mathbf x}=-\Sop\mathbf x.}
\label{eq:whitenedDynamics}
\end{equation}

\begin{mathteach}
Whitening is a change of metric, not a biological transformation of the organism. It expresses displacement in units of standing additive variation. Once genetic covariance has been made spherical, the joint effect of variance and curvature is represented by a symmetric matrix.
\end{mathteach}

\section{Proof that \(\Sop\) is symmetric positive definite}

Transpose \(\Sop\):
\begin{align}
\Sop^\top
&=(\G^{1/2}\Hh\G^{1/2})^\top\\
&=(\G^{1/2})^\top\Hh^\top(\G^{1/2})^\top\\
&=\G^{1/2}\Hh\G^{1/2}\\
&=\Sop.
\end{align}
The third line uses symmetry of both \(\G^{1/2}\) and \(\Hh\).

For any nonzero \(\mathbf x\),
\begin{align}
\mathbf x^\top\Sop\mathbf x
&=\mathbf x^\top\G^{1/2}\Hh\G^{1/2}\mathbf x\\
&=(\G^{1/2}\mathbf x)^\top\Hh(\G^{1/2}\mathbf x)\\
&>0,
\end{align}
because \(\G^{1/2}\mathbf x\neq\Zero\) and \(\Hh\succ0\). Hence \(\Sop\succ0\).

\begin{summarybox}
The operator \(\Sop\) is symmetric, has positive real eigenvalues, and admits an orthonormal eigenbasis. This is the precise reason it supports principal angles, stable subspaces, projector geometry, and symmetric perturbation theory.
\end{summarybox}

\section{Equivalence to the original dynamics}

The symmetric operator does not change the eigenvalues of the original response matrix. Observe that
\begin{align}
\G\Hh
&=\G^{1/2}(\G^{1/2}\Hh\G^{1/2})\G^{-1/2}\\
&=\G^{1/2}\Sop\G^{-1/2}.
\label{eq:similarityGH}
\end{align}
Thus \(\G\Hh\) and \(\Sop\) are similar matrices and have the same eigenvalues.

Suppose
\begin{equation}
\Sop\mathbf v_j=\lambda_j\mathbf v_j.
\end{equation}
Define the corresponding phenotype-space vector
\begin{equation}
\mathbf u_j=\G^{1/2}\mathbf v_j.
\label{eq:phenotypeChannelDirection}
\end{equation}
Then
\begin{align}
\G\Hh\mathbf u_j
&=\G\Hh\G^{1/2}\mathbf v_j\\
&=\G^{1/2}\Sop\mathbf v_j\\
&=\lambda_j\G^{1/2}\mathbf v_j\\
&=\lambda_j\mathbf u_j.
\end{align}
Therefore \(\mathbf u_j\) is a right eigenvector of \(\G\Hh\).

The same result can be written as a generalized symmetric eigenproblem. Starting from
\begin{equation}
\G\Hh\mathbf u_j=\lambda_j\mathbf u_j,
\end{equation}
left-multiply by \(\G^{-1}\):
\begin{equation}
\Hh\mathbf u_j=\lambda_j\G^{-1}\mathbf u_j.
\label{eq:generalisedEigenproblem}
\end{equation}

\inwords A response direction balances fitness curvature against inverse genetic covariance. It is an eigenvector of curvature measured in the genetic metric.

The phenotype-space directions are orthonormal under the \(\G^{-1}\) metric:
\begin{align}
\mathbf u_i^\top\G^{-1}\mathbf u_j
&=(\G^{1/2}\mathbf v_i)^\top\G^{-1}(\G^{1/2}\mathbf v_j)\\
&=\mathbf v_i^\top\G^{1/2}\G^{-1}\G^{1/2}\mathbf v_j\\
&=\mathbf v_i^\top\mathbf v_j\\
&=\delta_{ij}.
\label{eq:GinverseOrthogonality}
\end{align}

\begin{geombox}
In whitened space, channels are ordinary orthogonal axes. In phenotype space, the same directions are generally not Euclidean-orthogonal, but they are orthogonal when distances are measured in units of genetic covariance. The two descriptions encode the same dynamics.
\end{geombox}

\section{Solving the local dynamics mode by mode}

Let
\begin{equation}
\Sop=\Vv\boldsymbol\Lambda\Vv^\top,
\qquad
\boldsymbol\Lambda=\diag(\lambda_1,\ldots,\lambda_n),
\end{equation}
with \(\lambda_1\ge\cdots\ge\lambda_n>0\). Expand the initial displacement as
\begin{equation}
\mathbf x(0)=\sum_{j=1}^nc_j\mathbf v_j,
\qquad
c_j=\mathbf v_j^\top\mathbf x(0).
\end{equation}
Seek a solution of the form
\begin{equation}
\mathbf x(t)=\sum_{j=1}^na_j(t)\mathbf v_j.
\end{equation}
Substitute into \cref{eq:whitenedDynamics}:
\begin{align}
\sum_j\dot a_j(t)\mathbf v_j
&=-\Sop\sum_ja_j(t)\mathbf v_j\\
&=-\sum_ja_j(t)\lambda_j\mathbf v_j.
\end{align}
Orthogonality implies that each coefficient satisfies
\begin{equation}
\dot a_j=-\lambda_ja_j.
\end{equation}
The scalar solution is
\begin{equation}
a_j(t)=c_je^{-\lambda_jt}.
\end{equation}
Hence
\begin{equation}
\boxed{\mathbf x(t)=\sum_{j=1}^nc_je^{-\lambda_jt}\mathbf v_j
=e^{-\Sop t}\mathbf x(0).}
\label{eq:modeSolution}
\end{equation}

\inwords Every channel mode decays independently. A larger eigenvalue means faster local relaxation of displacement along that mode.

\begin{biobox}
The leading channel is the trait combination for which the product of available genetic variation and local restoring curvature generates the fastest local response. It is not necessarily the shallowest fitness ridge and not necessarily \(\mathbf g_{\max}\); it is the dominant mode of the combined system.
\end{biobox}

\begin{cautionbox}[title={A terminology condition in the ridge model}]
Later we model a ``channel ridge'' with high genetic variance \(g_{\parallel}\) and low curvature \(h_{\parallel}\) along a subspace. For that subspace to be the \emph{leading response channel}, it must also satisfy
\[
g_{\parallel}h_{\parallel}>g_{\perp}h_{\perp}.
\]
High variance along a shallow ridge does not automatically imply the largest response rate; the product matters.
\end{cautionbox}

\section{Recovery of \(\mathbf g_{\max}\)}

If curvature is isotropic,
\begin{equation}
\Hh=h\Ii,
\end{equation}
then
\begin{align}
\Sop
&=\G^{1/2}(h\Ii)\G^{1/2}\\
&=h\G.
\end{align}
Thus \(\Sop\) and \(\G\) have identical eigenvectors, and the leading channel is \(\mathbf g_{\max}\).

\begin{connectionbox}
This is the exact Schluter limit. The new framework does not reject genetic lines of least resistance. It identifies the condition under which they are complete: local fitness curvature must be isotropic, or at least commute with \(\G\) so that both matrices share axes.
\end{connectionbox}

\section{What changes if \(\G\) itself changes through time?}

If \(\G=\G(t)\), then
\begin{equation}
\mathbf x(t)=\G(t)^{-1/2}\mathbf d(t)
\end{equation}
and the product rule gives
\begin{equation}
\dot{\mathbf x}
=\frac{\dd}{\dd t}\left(\G^{-1/2}\right)\mathbf d
+\G^{-1/2}\dot{\mathbf d}.
\label{eq:timeVaryingWhitening}
\end{equation}
The second term still contains \(-\Sop\mathbf x\), but the first term records motion of the coordinate metric itself.

\begin{cautionbox}
Therefore \(\dot{\mathbf x}=-\Sop\mathbf x\) is an exact dynamical transformation when \(\G\) is fixed over the interval. When \(\G\) evolves, \(\Sop(t)\) remains the correct instantaneous symmetric response operator and its eigenspaces remain meaningful, but the full whitened trajectory includes an additional moving-metric term. Later we differentiate \(\Sop(t)\) directly to study channel rotation.
\end{cautionbox}

\chapter{A complete two-trait worked example}

\section{Write the operator in the eigenbasis of \(\G\)}

Choose coordinates aligned with the eigenvectors of \(\G\), so
\begin{equation}
\G=
\begin{bmatrix}
g_1&0\\0&g_2
\end{bmatrix},
\qquad g_1\ge g_2>0.
\end{equation}
Let
\begin{equation}
\Hh=
\begin{bmatrix}
h_{11}&h_{12}\\h_{12}&h_{22}
\end{bmatrix}.
\end{equation}
Then
\begin{equation}
\G^{1/2}=
\begin{bmatrix}
\sqrt{g_1}&0\\0&\sqrt{g_2}
\end{bmatrix}.
\end{equation}
Multiply in two steps:
\begin{align}
\G^{1/2}\Hh
&=
\begin{bmatrix}
\sqrt{g_1}h_{11}&\sqrt{g_1}h_{12}\\
\sqrt{g_2}h_{12}&\sqrt{g_2}h_{22}
\end{bmatrix},\\[0.5em]
\Sop=\G^{1/2}\Hh\G^{1/2}
&=
\begin{bmatrix}
g_1h_{11}&\sqrt{g_1g_2}h_{12}\\
\sqrt{g_1g_2}h_{12}&g_2h_{22}
\end{bmatrix}.
\label{eq:S2D}
\end{align}

\inwords The diagonal entries combine variance and curvature along the genetic axes. The off-diagonal entry is curvature that mixes the two genetic axes, scaled by their geometric-mean variance.

\section{Derive the eigenvector angle}

Write a generic symmetric two-by-two matrix as
\begin{equation}
\Sop=
\begin{bmatrix}
a&b\\b&c
\end{bmatrix}.
\end{equation}
A unit vector at angle \(\phi\) is
\begin{equation}
\mathbf v=
\begin{bmatrix}\cos\phi\\\sin\phi\end{bmatrix}.
\end{equation}
The Rayleigh quotient is
\begin{align}
r(\phi)
&=\mathbf v^\top\Sop\mathbf v\\
&=a\cos^2\phi+2b\sin\phi\cos\phi+c\sin^2\phi.
\end{align}
Use
\begin{equation}
\cos^2\phi=\frac{1+\cos2\phi}{2},
\quad
\sin^2\phi=\frac{1-\cos2\phi}{2},
\quad
2\sin\phi\cos\phi=\sin2\phi,
\end{equation}
to obtain
\begin{equation}
r(\phi)=\frac{a+c}{2}+\frac{a-c}{2}\cos2\phi+b\sin2\phi.
\end{equation}
Differentiate:
\begin{equation}
\frac{\dd r}{\dd\phi}=-(a-c)\sin2\phi+2b\cos2\phi.
\end{equation}
At an eigenvector direction, the Rayleigh quotient is stationary, so
\begin{equation}
-(a-c)\sin2\phi+2b\cos2\phi=0.
\end{equation}
Rearranging,
\begin{equation}
\tan2\phi=\frac{2b}{a-c}.
\end{equation}
The quadrant-safe expression is
\begin{equation}
\boxed{\phi=\half\operatorname{atan2}(2b,a-c).}
\label{eq:2DeigenAngle}
\end{equation}
Substituting the entries from \cref{eq:S2D},
\begin{equation}
\boxed{
\tan2\phi=
\frac{2\sqrt{g_1g_2}\,h_{12}}
{g_1h_{11}-g_2h_{22}}.}
\label{eq:2DchannelAngle}
\end{equation}

\begin{mathteach}
The numerator is rotational forcing: it mixes the two axes. The denominator is the separation between their uncoupled response rates. Rotation is large when mixing is strong or when the two rates are nearly tied.
\end{mathteach}

\begin{geombox}
The two-dimensional formula already contains the later perturbation theory. The off-diagonal term tries to turn the axes; the eigenvalue gap resists that turn. Near a tie, a small off-diagonal perturbation can cause a large angle.
\end{geombox}

\section{Map the whitened angle back to phenotype space}

The whitened eigenvector is
\begin{equation}
\mathbf v=
\begin{bmatrix}\cos\phi\\\sin\phi\end{bmatrix}.
\end{equation}
The phenotype-space response direction is
\begin{equation}
\mathbf u=\G^{1/2}\mathbf v
=
\begin{bmatrix}
\sqrt{g_1}\cos\phi\\
\sqrt{g_2}\sin\phi
\end{bmatrix}.
\end{equation}
If \(\alpha\) is its ordinary Euclidean angle,
\begin{equation}
\tan\alpha
=\frac{u_2}{u_1}
=\sqrt{\frac{g_2}{g_1}}\tan\phi.
\label{eq:phenotypeAngle2D}
\end{equation}

\inwords Whitening and phenotype space do not assign the same Euclidean angle when genetic variance is anisotropic. The phenotype direction is stretched more along high-variance axes when mapped back from whitened coordinates.

\section{Three informative special cases}

\subsection{Isotropic curvature}
If \(\Hh=h\Ii\), then \(h_{12}=0\), \(h_{11}=h_{22}=h\), and \cref{eq:2DchannelAngle} gives \(\phi=0\) when \(g_1>g_2\). The leading channel is \(\mathbf g_{\max}\).

\subsection{Aligned but anisotropic curvature}
If \(h_{12}=0\), the matrices share axes. There is no smooth rotation. The leading direction is axis 1 when
\begin{equation}
g_1h_{11}>g_2h_{22}
\end{equation}
and axis 2 when the inequality reverses. At
\begin{equation}
g_1h_{11}=g_2h_{22},
\end{equation}
the eigenvalues are equal and the rank-one leading direction is undefined.

\begin{biobox}
Selection can change the strength of response modes without rotating them. Rotation requires a change in trait combinations, represented locally by cross-channel terms. A population can therefore experience strong selection and major eigenvalue change while its channel orientation remains constant.
\end{biobox}

\subsection{Misaligned curvature}
If \(h_{12}\neq0\), curvature couples the genetic axes and the leading channel rotates smoothly. The rotation becomes most sensitive as \(g_1h_{11}-g_2h_{22}\) approaches zero.

\begin{figure}[H]
\centering
\includegraphics[width=0.72\textwidth]{figures/figS1_channel_rotation_checked.pdf}
\caption{Two-trait illustration of rotation away from \(\mathbf g_{\max}\) as curvature anisotropy increases. Genetic correlation changes the orientation and susceptibility of the joint operator.}
\label{fig:2Drotation}
\end{figure}

\chapter{Channels as invariant subspaces}

\section{Why a channel need not be one-dimensional}

Let the eigendecomposition of \(\Sop\) be
\begin{equation}
\Sop=\sum_{j=1}^n\lambda_j\mathbf v_j\mathbf v_j^\top,
\qquad
\lambda_1\ge\cdots\ge\lambda_n.
\end{equation}
A rank-\(k\) channel basis is
\begin{equation}
\Vv_k=
\begin{bmatrix}
\mathbf v_1&\cdots&\mathbf v_k
\end{bmatrix}.
\end{equation}
The channel itself is the subspace
\begin{equation}
\mathcal C_k=\Span(\Vv_k).
\end{equation}

\inwords A basis is a particular set of coordinate vectors. A subspace is the collection of all linear combinations of those vectors. Rotating the basis within the same subspace changes coordinates but not the biological channel.

\section{The projector removes arbitrary basis choice}

Define the orthogonal projector
\begin{equation}
\Pp_k=\Vv_k\Vv_k^\top.
\label{eq:projectorDefinition}
\end{equation}
For any vector \(\mathbf x\),
\begin{equation}
\Pp_k\mathbf x
\end{equation}
is its projection onto the channel.

Suppose \(\widetilde{\Vv}_k=\Vv_k\mathbf O\), where \(\mathbf O\) is any \(k\times k\) orthogonal rotation. Then
\begin{align}
\widetilde{\Vv}_k\widetilde{\Vv}_k^\top
&=\Vv_k\mathbf O\mathbf O^\top\Vv_k^\top\\
&=\Vv_k\Vv_k^\top\\
&=\Pp_k.
\end{align}
Thus the projector depends only on the subspace, not on the arbitrary basis used to represent it.

\begin{mathteach}
The set of all \(k\)-dimensional subspaces of \(\R^n\) is the Grassmann manifold \(\Grass(k,n)\). Channel evolution is therefore naturally a trajectory on a Grassmannian, not merely a trajectory of individual eigenvectors.
\end{mathteach}

\section{The biologically relevant spectral gap}

The gap separating the selected channel from excluded modes is
\begin{equation}
\delta_k=\lambda_k-\lambda_{k+1}.
\label{eq:kGap}
\end{equation}

If \(\lambda_1\approx\lambda_2\) but \(\lambda_2\gg\lambda_3\), the leading vector is unstable but the leading two-dimensional subspace is stable. A forced rank-one analysis would report dramatic ``rotation'' that is largely an arbitrary exchange between two nearly tied modes.

\begin{biobox}
The appropriate channel dimension should reflect the number of response modes that are biologically distinguishable from the remainder. A channel boundary is always relative to \(k\). There is no universal reason to define it only by \(\lambda_1-\lambda_2\).
\end{biobox}

\section{Choosing \(k\) in practice}

A defensible choice can combine:
\begin{enumerate}[leftmargin=1.6em]
\item a posterior distribution for eigengaps;
\item stability of the subspace under resampling;
\item cumulative response rate \(\sum_{j\le k}\lambda_j/\sum_j\lambda_j\);
\item predictive performance for observed response or hybrid fitness;
\item a preregistered biological hypothesis about modular trait combinations.
\end{enumerate}

\begin{cautionbox}
Choosing \(k\) after inspecting which dimension gives the strongest association with reproductive isolation will inflate evidence. The dimension should be selected by independent stability or predictive criteria and propagated through uncertainty.
\end{cautionbox}

\part{Comparing channels and deriving the geometric snowball}

\chapter{Principal angles, projectors, and overlap between subspaces}

\section{The one-dimensional case}

For two unit vectors \(\mathbf u_1\) and \(\mathbf u_2\), the acute angle is
\begin{equation}
\psi=\arccos\left(|\mathbf u_1^\top\mathbf u_2|\right).
\end{equation}
The absolute value appears because a line has no inherent sign: \(\mathbf u\) and \(-\mathbf u\) span the same one-dimensional subspace.

\inwords The inner product is the length of one unit vector retained when it is projected onto the other. Perfect alignment gives one; orthogonality gives zero.

\section{Principal angles for \(k\)-dimensional channels}

Let \(\Vv_1,\Vv_2\in\R^{n\times k}\) contain orthonormal bases:
\begin{equation}
\Vv_i^\top\Vv_i=\Ii_k.
\end{equation}
Form the cross-basis matrix
\begin{equation}
\mathbf C=\Vv_1^\top\Vv_2.
\end{equation}
Take its singular value decomposition:
\begin{equation}
\mathbf C=\mathbf L\diag(\sigma_1,\ldots,\sigma_k)\mathbf R^\top,
\qquad
1\ge\sigma_1\ge\cdots\ge\sigma_k\ge0.
\end{equation}
The principal angles are defined by
\begin{equation}
\boxed{\cos\psi_j=\sigma_j.}
\label{eq:principalAngles}
\end{equation}

\begin{mathteach}
The matrices \(\mathbf L\) and \(\mathbf R\) merely choose the paired directions within the two subspaces that align best, second-best, and so on. The singular values are invariant to arbitrary rotations of either basis and therefore belong to the subspaces themselves.
\end{mathteach}

\section{Projector form of chordal distance}

Let
\begin{equation}
\Pp_i=\Vv_i\Vv_i^\top.
\end{equation}
First compute
\begin{align}
\tr(\Pp_1\Pp_2)
&=\tr(\Vv_1\Vv_1^\top\Vv_2\Vv_2^\top)\\
&=\tr(\Vv_1^\top\Vv_2\Vv_2^\top\Vv_1)\\
&=\tr(\mathbf C\mathbf C^\top)\\
&=\normF{\mathbf C}^2\\
&=\sum_{j=1}^k\sigma_j^2\\
&=\sum_{j=1}^k\cos^2\psi_j.
\label{eq:traceProjectorOverlap}
\end{align}
The second line uses cyclic invariance of trace. The fifth uses the fact that the squared Frobenius norm is the sum of squared singular values.

Now expand the projector difference:
\begin{align}
\normF{\Pp_1-\Pp_2}^2
&=\tr\left[(\Pp_1-\Pp_2)^2\right]\\
&=\tr(\Pp_1^2)+\tr(\Pp_2^2)-2\tr(\Pp_1\Pp_2).
\end{align}
Because projectors satisfy \(\Pp_i^2=\Pp_i\) and \(\tr(\Pp_i)=k\),
\begin{align}
\normF{\Pp_1-\Pp_2}^2
&=2k-2\tr(\Pp_1\Pp_2)\\
&=2\sum_{j=1}^k(1-\cos^2\psi_j)\\
&=2\sum_{j=1}^k\sin^2\psi_j.
\end{align}
Therefore the squared chordal distance is
\begin{equation}
\boxed{
d_C^2(\mathcal C_1,\mathcal C_2)
=\half\normF{\Pp_1-\Pp_2}^2
=\sum_{j=1}^k\sin^2\psi_j.}
\label{eq:chordalIdentity}
\end{equation}

\inwords Chordal distance is the total squared amount by which the paired channel directions fail to project onto one another.

\begin{geombox}
The projector map embeds the curved Grassmann manifold into the flat Euclidean space of symmetric matrices. Chordal distance is ordinary Euclidean distance between the embedded projectors, scaled by one half.
\end{geombox}

\section{Volume overlap}

From the SVD,
\begin{align}
\det(\Vv_1^\top\Vv_2)^2
&=\det(\mathbf L)^2
\det[\diag(\sigma_j)]^2
\det(\mathbf R)^2\\
&=\prod_{j=1}^k\sigma_j^2\\
&=\prod_{j=1}^k\cos^2\psi_j,
\end{align}
because orthogonal matrices have determinant \(\pm1\).

Define the bounded channel-overlap deficit
\begin{equation}
\boxed{
I_V=1-\det(\Vv_1^\top\Vv_2)^2
=1-\prod_{j=1}^k\cos^2\psi_j.}
\label{eq:volumeDeficit}
\end{equation}

\inwords The product \(\prod_j\cos\psi_j\) is the factor by which projecting a unit \(k\)-dimensional volume element from one channel into the other shrinks that volume. Squaring removes orientation and gives a convenient overlap measure.

A related additive quantity is the log-overlap loss
\begin{equation}
L_V=-\log(1-I_V)
=-2\sum_{j=1}^k\log(\cos\psi_j).
\label{eq:logOverlap}
\end{equation}
It is additive across principal directions, but it is still a geometric descriptor rather than a calibrated fitness load.

\section{Small-angle expansion}

For small \(\psi\),
\begin{align}
\sin^2\psi
&=\psi^2-\frac{\psi^4}{3}+O(\psi^6),\\
\cos^2\psi
&=1-\psi^2+\frac{\psi^4}{3}+O(\psi^6).
\end{align}
Therefore
\begin{equation}
d_C^2
=\sum_j\psi_j^2-\frac13\sum_j\psi_j^4+O(\normtwo{\psiv}^6).
\label{eq:chordalExpansion}
\end{equation}
For the volume deficit,
\begin{align}
\prod_j\cos^2\psi_j
&=\prod_j\left(1-\psi_j^2+\frac{\psi_j^4}{3}+O(\psi_j^6)\right)\\
&=1-\sum_j\psi_j^2
+\frac13\sum_j\psi_j^4
+\sum_{i<j}\psi_i^2\psi_j^2
+O(\normtwo{\psiv}^6),
\end{align}
so
\begin{equation}
I_V
=\sum_j\psi_j^2
-\frac13\sum_j\psi_j^4
-\sum_{i<j}\psi_i^2\psi_j^2
+O(\normtwo{\psiv}^6).
\label{eq:volumeExpansion}
\end{equation}
Thus
\begin{equation}
\boxed{I_V=d_C^2+O(\normtwo{\psiv}^4)
=\sum_j\psi_j^2+O(\normtwo{\psiv}^4).}
\label{eq:overlapLocalEquivalence}
\end{equation}

\begin{connectionbox}
The quadratic snowball does not depend on choosing the determinant deficit rather than chordal distance. All smooth overlap deficits with a nonzero quadratic term agree locally up to scaling. The choice matters at larger angles and for interpretation, but not for the early exponent.
\end{connectionbox}

\begin{figure}[H]
\centering
\includegraphics[width=0.96\textwidth]{figures/pedagogy_principal_angles.pdf}
\caption{Three equivalent geometric readings of channel alignment: angle, projection length, and projected \(k\)-volume.}
\label{fig:principalAnglePedagogy}
\end{figure}

\chapter{Comparing channels when populations have different \(\G\) matrices}

\section{Why population-specific whitened spaces cannot be compared naively}

Population \(i\) has a whitened coordinate
\begin{equation}
\mathbf x_i=\G_i^{-1/2}\mathbf d.
\end{equation}
If \(\G_1\neq\G_2\), the same phenotype displacement receives different whitened coordinates in the two populations. Therefore a Euclidean angle between eigenvectors of \(\Sop_1\) and \(\Sop_2\) compares vectors that live in different metrics.

\begin{cautionbox}
This is not a minor numerical detail. A raw angle between population-specific whitened eigenvectors can change when one trait is converted from centimetres to millimetres. A biologically meaningful comparison must use one shared metric.
\end{cautionbox}

\section{Map channels back into phenotype space}

Let \(\Vv_i\in\R^{n\times k}\) be the orthonormal channel basis of \(\Sop_i\). Define
\begin{equation}
\mathbf U_i=\G_i^{1/2}\Vv_i.
\label{eq:mapChannelsPhenotype}
\end{equation}
The columns of \(\mathbf U_i\) span the corresponding phenotype-space response subspace.

Choose a common SPD metric \(\M\). Possible choices include ancestral \(\G_0^{-1}\), pooled \(\G^{-1}\), or a preregistered biologically motivated measurement metric.

\begin{cautionbox}
Coordinate invariance means that the result is unchanged when the same biological metric is represented in different units or coordinates. It does \emph{not} mean that every scientifically plausible choice of \(\M\) gives the same angle. The metric defines what counts as one unit of phenotypic displacement, so its choice is part of the biological model and should be stated and, where possible, tested for sensitivity.
\end{cautionbox}

The Gram matrix of \(\mathbf U_i\) under \(\M\) is
\begin{equation}
\mathbf A_i=\mathbf U_i^\top\M\mathbf U_i.
\end{equation}
Because \(\mathbf U_i\) has full column rank and \(\M\succ0\), \(\mathbf A_i\succ0\).

Define the \(\M\)-orthonormal basis
\begin{equation}
\boxed{
\Qq_i=\mathbf U_i\mathbf A_i^{-1/2}
=\mathbf U_i(\mathbf U_i^\top\M\mathbf U_i)^{-1/2}.}
\label{eq:MOrthonormalBasis}
\end{equation}
Check the normalization:
\begin{align}
\Qq_i^\top\M\Qq_i
&=\mathbf A_i^{-1/2}\mathbf U_i^\top\M\mathbf U_i\mathbf A_i^{-1/2}\\
&=\mathbf A_i^{-1/2}\mathbf A_i\mathbf A_i^{-1/2}\\
&=\Ii_k.
\end{align}

The common-metric principal-angle cosines are the singular values of
\begin{equation}
\boxed{\Qq_1^\top\M\Qq_2.}
\label{eq:commonMetricCrossMatrix}
\end{equation}

\inwords Each channel is first expressed in the same phenotype space, then normalized using the same rule for length and angle. Only after that are the subspaces compared.

\section{Proof of coordinate invariance}

Under \(\mathbf z'=\mathbf T\mathbf z\), phenotype-space channel bases transform as
\begin{equation}
\mathbf U_i'=\mathbf T\mathbf U_i.
\end{equation}
The metric transforms as
\begin{equation}
\M'=\mathbf T^{-\top}\M\mathbf T^{-1}.
\end{equation}
Then
\begin{align}
(\mathbf U_i')^\top\M'\mathbf U_i'
&=(\mathbf T\mathbf U_i)^\top
(\mathbf T^{-\top}\M\mathbf T^{-1})
(\mathbf T\mathbf U_i)\\
&=\mathbf U_i^\top\mathbf T^\top\mathbf T^{-\top}
\M
\mathbf T^{-1}\mathbf T\mathbf U_i\\
&=\mathbf U_i^\top\M\mathbf U_i\\
&=\mathbf A_i.
\end{align}
Therefore
\begin{equation}
\Qq_i'=\mathbf T\mathbf U_i\mathbf A_i^{-1/2}
=\mathbf T\Qq_i.
\end{equation}
The cross matrix becomes
\begin{align}
(\Qq_1')^\top\M'\Qq_2'
&=(\mathbf T\Qq_1)^\top
(\mathbf T^{-\top}\M\mathbf T^{-1})
(\mathbf T\Qq_2)\\
&=\Qq_1^\top\M\Qq_2.
\end{align}
It is exactly unchanged. Consequently all singular values and principal angles are coordinate invariant.

\begin{summarybox}
When \(\G_1\neq\G_2\), the invariant comparison is: map each channel to phenotype space, orthonormalize both under one common metric, and compute principal angles from \(\Qq_1^\top\M\Qq_2\).
\end{summarybox}

\section{Recovery of the shared-\(\G\) comparison}

If \(\G_1=\G_2=\G\) and choose \(\M=\G^{-1}\), then
\begin{align}
\mathbf U_i^\top\M\mathbf U_i
&=(\G^{1/2}\Vv_i)^\top\G^{-1}(\G^{1/2}\Vv_i)\\
&=\Vv_i^\top\Vv_i\\
&=\Ii_k.
\end{align}
Thus \(\Qq_i=\mathbf U_i\), and
\begin{align}
\Qq_1^\top\M\Qq_2
&=\Vv_1^\top\Vv_2.
\end{align}
The common-metric method therefore reduces exactly to the ordinary whitened-space comparison when a common \(\G\) exists.

\chapter{First-order eigenspace perturbation}

\section{Set up two smoothly diverging populations}

\begin{cautionbox}
The theorem first derived in this chapter is an abstract perturbation result for symmetric operators represented in one fixed Euclidean inner-product space. It applies directly to biological channels when the whitening transformation is shared and fixed---for example, when both populations are compared using the same ancestral \(\G\). When \(\G_i(t)\) itself evolves, the phenotype-space channel also changes through the map \(\G_i^{1/2}\). The fully coordinate-invariant coefficient including that additional term is derived later in this chapter.
\end{cautionbox}

Assume both populations begin from the same symmetric operator \(\Sop_0\) and then diverge smoothly:
\begin{equation}
\Sop_i(t)=\Sop_0+t\Aa_i+O(t^2),
\qquad i\in\{1,2\}.
\label{eq:operatorPaths}
\end{equation}
The matrix \(\Aa_i=\dot\Sop_i(0)\) is the initial operator velocity.

Let
\begin{equation}
\Sop_0\mathbf v_j=\lambda_j\mathbf v_j,
\qquad
\lambda_1\ge\cdots\ge\lambda_n.
\end{equation}
For a rank-\(k\) channel, collect
\begin{equation}
\Vv=
\begin{bmatrix}\mathbf v_1&\cdots&\mathbf v_k\end{bmatrix},
\qquad
\Ww=
\begin{bmatrix}\mathbf w_1&\cdots&\mathbf w_{n-k}\end{bmatrix},
\end{equation}
where \(\mathbf w_a=\mathbf v_{k+a}\) span the orthogonal complement.

\section{Derive the change of a single eigenvector}

For one population, suppress the subscript \(i\). Let
\begin{equation}
\Sop(t)\mathbf v_b(t)=\lambda_b(t)\mathbf v_b(t).
\label{eq:eigenEquationTime}
\end{equation}
Differentiate at \(t=0\):
\begin{equation}
\Aa\mathbf v_b+\Sop_0\dot{\mathbf v}_b
=\dot\lambda_b\mathbf v_b+\lambda_b\dot{\mathbf v}_b.
\label{eq:eigenDerivativeRaw}
\end{equation}

The eigenvector is normalized:
\begin{equation}
\mathbf v_b(t)^\top\mathbf v_b(t)=1.
\end{equation}
Differentiate:
\begin{equation}
2\mathbf v_b^\top\dot{\mathbf v}_b=0,
\end{equation}
so
\begin{equation}
\mathbf v_b^\top\dot{\mathbf v}_b=0.
\label{eq:eigenvectorDerivativeOrthogonal}
\end{equation}
Thus the derivative has no component along the eigenvector itself.

Left-multiply \cref{eq:eigenDerivativeRaw} by \(\mathbf v_b^\top\):
\begin{align}
\mathbf v_b^\top\Aa\mathbf v_b
+\mathbf v_b^\top\Sop_0\dot{\mathbf v}_b
&=\dot\lambda_b
+\lambda_b\mathbf v_b^\top\dot{\mathbf v}_b.
\end{align}
Because \(\mathbf v_b^\top\Sop_0=\lambda_b\mathbf v_b^\top\) and \cref{eq:eigenvectorDerivativeOrthogonal} holds, the terms involving \(\dot{\mathbf v}_b\) cancel:
\begin{equation}
\boxed{\dot\lambda_b=\mathbf v_b^\top\Aa\mathbf v_b.}
\label{eq:eigenvalueDerivative}
\end{equation}

\inwords The first-order change in an eigenvalue is the perturbation measured along its own eigenvector.

Now left-multiply \cref{eq:eigenDerivativeRaw} by a different eigenvector \(\mathbf v_a^\top\), with \(a\neq b\):
\begin{align}
\mathbf v_a^\top\Aa\mathbf v_b
+\lambda_a\mathbf v_a^\top\dot{\mathbf v}_b
&=\lambda_b\mathbf v_a^\top\dot{\mathbf v}_b,
\end{align}
because \(\mathbf v_a^\top\mathbf v_b=0\). Rearranging,
\begin{equation}
\boxed{
\mathbf v_a^\top\dot{\mathbf v}_b
=\frac{\mathbf v_a^\top\Aa\mathbf v_b}{\lambda_b-\lambda_a}.}
\label{eq:eigenvectorDerivativeCoefficient}
\end{equation}
Therefore
\begin{equation}
\dot{\mathbf v}_b
=\sum_{a\neq b}
\frac{\mathbf v_a^\top\Aa\mathbf v_b}{\lambda_b-\lambda_a}\mathbf v_a,
\label{eq:fullEigenvectorDerivative}
\end{equation}
up to a sign convention or gauge inside exactly degenerate eigenspaces.

\begin{mathteach}
The numerator is off-diagonal forcing: it asks whether the perturbation maps the current eigenvector into another eigendirection. The denominator is resistance supplied by eigenvalue separation. A large gap stabilizes orientation; a small gap amplifies rotation.
\end{mathteach}

\section{From eigenvectors to a rank-\(k\) subspace}

Only components that leave the channel subspace change the subspace itself. Rotations among the first \(k\) vectors merely change the internal basis. The complement component is
\begin{equation}
\Ww^\top\dot{\Vv}=\cK(\Aa),
\end{equation}
with entries
\begin{equation}
\boxed{
\cK(\Aa)_{ab}
=\frac{\mathbf w_a^\top\Aa\mathbf v_b}
{\lambda_b-\lambda_{k+a}}.}
\label{eq:tangentMap}
\end{equation}
We call \(\cT:\Aa\mapsto\cK(\Aa)\) the local tangent map from operator change to channel rotation.

For the two populations, the relative tangent forcing is
\begin{equation}
\boxed{
\cK=\cK(\Aa_2-\Aa_1),
\qquad
\cK_{ab}=
\frac{\mathbf w_a^\top(\Aa_2-\Aa_1)\mathbf v_b}
{\lambda_b-\lambda_{k+a}}.}
\label{eq:relativeK}
\end{equation}

\section{Derive the local quadratic channel law using projectors}

Choose the parallel-transport gauge \(\Vv^\top\dot{\Vv}_i=\Zero\). Then
\begin{equation}
\dot{\Vv}_i=\Ww\cK_i,
\end{equation}
where \(\cK_i=\cK(\Aa_i)\). Thus
\begin{equation}
\Vv_i(t)=\Vv+t\Ww\cK_i+O(t^2).
\end{equation}
The projector is
\begin{align}
\Pp_i(t)
&=\Vv_i(t)\Vv_i(t)^\top\\
&=\left(\Vv+t\Ww\cK_i+O(t^2)\right)
\left(\Vv^\top+t\cK_i^\top\Ww^\top+O(t^2)\right)\\
&=\Vv\Vv^\top
+t\Ww\cK_i\Vv^\top
+t\Vv\cK_i^\top\Ww^\top
+O(t^2).
\end{align}
Subtract the two populations:
\begin{equation}
\Pp_2(t)-\Pp_1(t)
=t\left(\Ww\cK\Vv^\top+\Vv\cK^\top\Ww^\top\right)+O(t^2).
\label{eq:projectorDifferenceExpansion}
\end{equation}

Let
\begin{equation}
\mathbf B=\Ww\cK\Vv^\top.
\end{equation}
Then \(\mathbf B^\top=\Vv\cK^\top\Ww^\top\). Because \(\Vv^\top\Ww=\Zero\),
\begin{equation}
\innerF{\mathbf B}{\mathbf B^\top}
=\tr(\mathbf B^2)=0.
\end{equation}
Also
\begin{align}
\normF{\mathbf B}^2
&=\tr(\Vv\cK^\top\Ww^\top\Ww\cK\Vv^\top)\\
&=\tr(\Vv\cK^\top\cK\Vv^\top)\\
&=\tr(\cK^\top\cK)\\
&=\normF{\cK}^2.
\end{align}
Therefore
\begin{align}
\normF{\Pp_2(t)-\Pp_1(t)}^2
&=t^2\normF{\mathbf B+\mathbf B^\top}^2+O(t^3)\\
&=t^2\left(\normF{\mathbf B}^2+\normF{\mathbf B^\top}^2
+2\innerF{\mathbf B}{\mathbf B^\top}\right)+O(t^3)\\
&=2t^2\normF{\cK}^2+O(t^3).
\end{align}
Using \cref{eq:chordalIdentity},
\begin{theorem}[Local geometric snowball]
If the rank-\(k\) channel of \(\Sop_0\) is separated from its complement by a nonzero gap and the two operators diverge smoothly as in \cref{eq:operatorPaths}, then
\begin{equation}
\boxed{
d_C^2(t)=\sum_{j=1}^k\sin^2\psi_j(t)
=t^2\normF{\cK}^2+O(t^3).}
\label{eq:localSnowballTheorem}
\end{equation}
\end{theorem}

\inwords Channel angle grows linearly at first, while any squared overlap deficit grows quadratically. The coefficient is not an unspecified constant: it is the sum of squared cross-channel forcings, each divided by the relevant squared eigengap.

For \(k=1\),
\begin{equation}
\boxed{
\theta_{12}^2(t)
=t^2\sum_{a=2}^n
\left[
\frac{\mathbf v_a^\top(\Aa_2-\Aa_1)\mathbf v_1}
{\lambda_1-\lambda_a}
\right]^2+O(t^3).}
\label{eq:rankOneSnowball}
\end{equation}

\begin{summarybox}
In a fixed shared representation, the exact local susceptibility is
\[
\chi_k^2=\normF{\cK}^2.
\]
A small gap is only one ingredient. Rapid rotation also requires the ecological or genetic perturbation to couple the selected channel to excluded directions. When \(\G\) evolves, the fully invariant susceptibility is \(\chi_M^2=\normF{\mathbf K_M}^2\), derived in the next section.
\end{summarybox}

\begin{figure}[H]
\centering
\includegraphics[width=0.84\textwidth]{figures/01_local_quadratic_law.pdf}
\caption{Numerical verification of the local coefficient in \cref{eq:localSnowballTheorem}. The approximation is exact asymptotically; higher-order curvature of the path becomes visible at larger divergence.}
\label{fig:localQuadraticCheck}
\end{figure}

\section{The fully invariant local coefficient when \(\G\) evolves}

The preceding theorem describes the rotation of the leading eigenspace of \(\Sop_i(t)\) inside one fixed Euclidean representation. We now derive the coefficient for the biologically relevant phenotype-space channels when \(\G_i(t)\) itself changes.

Assume the two populations share the same starting point,
\begin{equation}
\G_i(0)=\G_0,
\qquad
\Hh_i(0)=\Hh_0,
\qquad
\Sop_0=\mathbf R_0\Hh_0\mathbf R_0,
\qquad
\mathbf R_0=\G_0^{1/2}.
\end{equation}
Let \(\Vv\) and \(\Ww\) contain the starting channel and complementary eigenvectors of \(\Sop_0\). The phenotype-space channel basis for population \(i\) is
\begin{equation}
\mathbf U_i(t)=\mathbf R_i(t)\Vv_i(t),
\qquad
\mathbf R_i(t)=\G_i(t)^{1/2}.
\label{eq:timeVaryingPhenotypeChannel}
\end{equation}

Choose the fixed common metric \(\M\succ0\) introduced in the previous chapter. Multiplication by \(\M^{1/2}\) converts that metric into an ordinary Euclidean one. Define
\begin{equation}
\mathbf B_i(t)=\M^{1/2}\mathbf U_i(t),
\qquad
\boldsymbol\Gamma_i(t)=\mathbf B_i(t)^\top\mathbf B_i(t).
\label{eq:metricEuclideanization}
\end{equation}
Then an ordinary Euclidean orthonormal basis for the transformed channel is
\begin{equation}
\widehat{\Qq}_i(t)
=\mathbf B_i(t)\boldsymbol\Gamma_i(t)^{-1/2}.
\label{eq:EuclideanizedChannelBasis}
\end{equation}
At the common origin, write
\begin{equation}
\mathbf B_0=\M^{1/2}\mathbf R_0\Vv,
\qquad
\boldsymbol\Gamma_0=\mathbf B_0^\top\mathbf B_0,
\qquad
\widehat{\Qq}_0=\mathbf B_0\boldsymbol\Gamma_0^{-1/2}.
\end{equation}
Let \(\widehat{\Qq}_{\perp}\) be any Euclidean orthonormal basis for the complement of \(\widehat{\Qq}_0\):
\begin{equation}
\widehat{\Qq}_{\perp}^{\top}\widehat{\Qq}_0=\Zero.
\end{equation}

\subsection{Differentiate the metric-normalized basis}

Differentiate \cref{eq:EuclideanizedChannelBasis} at \(t=0\):
\begin{align}
\dot{\widehat{\Qq}}_i
&=\dot{\mathbf B}_i\boldsymbol\Gamma_0^{-1/2}
+\mathbf B_0
\left.\frac{\dd}{\dd t}\boldsymbol\Gamma_i(t)^{-1/2}\right|_{t=0}.
\label{eq:QhatDerivativeRaw}
\end{align}
The second term changes the normalization and the internal basis, but it lies entirely in the starting channel because it is left-multiplied by \(\mathbf B_0\). Projecting onto the complement removes it:
\begin{align}
\widehat{\Qq}_{\perp}^{\top}\dot{\widehat{\Qq}}_i
&=\widehat{\Qq}_{\perp}^{\top}\dot{\mathbf B}_i
\boldsymbol\Gamma_0^{-1/2}
+\underbrace{\widehat{\Qq}_{\perp}^{\top}\mathbf B_0}_{\Zero}
\left.\frac{\dd}{\dd t}\boldsymbol\Gamma_i^{-1/2}\right|_0\\
&=\widehat{\Qq}_{\perp}^{\top}\dot{\mathbf B}_i
\boldsymbol\Gamma_0^{-1/2}.
\label{eq:metricTangentProjection}
\end{align}

\inwords Only motion that leaves the current subspace changes the subspace. Renormalizing vectors or rotating the basis within the same channel does not contribute to a principal angle.

Define the population-specific common-metric tangent matrix
\begin{equation}
\mathbf K_{M,i}
=\widehat{\Qq}_{\perp}^{\top}\dot{\widehat{\Qq}}_i.
\label{eq:KMiDefinition}
\end{equation}
By \cref{eq:metricTangentProjection},
\begin{equation}
\mathbf K_{M,i}
=\widehat{\Qq}_{\perp}^{\top}\M^{1/2}\dot{\mathbf U}_i
\boldsymbol\Gamma_0^{-1/2}.
\label{eq:KMiFromUdot}
\end{equation}

\subsection{Separate map-back deformation from eigenspace rotation}

Differentiate \(\mathbf U_i=\mathbf R_i\Vv_i\):
\begin{equation}
\dot{\mathbf U}_i
=\dot{\mathbf R}_i\Vv+\mathbf R_0\dot{\Vv}_i.
\label{eq:UdotProductRule}
\end{equation}
Let
\begin{equation}
\mathbf X_i=\dot{\mathbf R}_i.
\end{equation}
Because \(\mathbf R_i^2=\G_i\), \(\mathbf X_i\) is the unique symmetric solution of
\begin{equation}
\boxed{
\mathbf R_0\mathbf X_i+\mathbf X_i\mathbf R_0=\dot\G_i.}
\label{eq:localSylvesterXi}
\end{equation}
The derivation and explicit solution of this Sylvester equation are given again in \cref{eq:SylvesterSqrt,eq:SqrtDerivativeElements}.

The derivative of the leading eigenspace of \(\Sop_i\) can be decomposed as
\begin{equation}
\dot{\Vv}_i
=\Ww\mathbf K_{S,i}+\Vv\boldsymbol\Omega_i,
\label{eq:VdotExternalInternal}
\end{equation}
where \(\boldsymbol\Omega_i\) is an internal basis rotation and
\begin{equation}
(\mathbf K_{S,i})_{ab}
=\frac{\mathbf w_a^\top\dot\Sop_i\mathbf v_b}
{\lambda_b-\lambda_{k+a}}.
\label{eq:KSiDefinition}
\end{equation}
Substitute \cref{eq:VdotExternalInternal} into \cref{eq:UdotProductRule}:
\begin{equation}
\dot{\mathbf U}_i
=\mathbf X_i\Vv
+\mathbf R_0\Ww\mathbf K_{S,i}
+\mathbf R_0\Vv\boldsymbol\Omega_i.
\end{equation}
The last term is internal to the phenotype-space channel. Indeed,
\begin{align}
\widehat{\Qq}_{\perp}^{\top}\M^{1/2}
\mathbf R_0\Vv\boldsymbol\Omega_i
&=\widehat{\Qq}_{\perp}^{\top}\mathbf B_0\boldsymbol\Omega_i\\
&=\Zero.
\end{align}
Hence
\begin{equation}
\boxed{
\mathbf K_{M,i}
=\widehat{\Qq}_{\perp}^{\top}\M^{1/2}
\left(\mathbf X_i\Vv+\mathbf R_0\Ww\mathbf K_{S,i}\right)
\boldsymbol\Gamma_0^{-1/2}.}
\label{eq:KMiFull}
\end{equation}

\begin{mathteach}
\Cref{eq:KMiFull} is a product-rule decomposition for a moving subspace. The term \(\mathbf X_i\Vv\) records how evolution of \(\G_i^{1/2}\) deforms the map from whitened space back to phenotype space. The term \(\mathbf R_0\Ww\mathbf K_{S,i}\) records rotation of the eigenspace of \(\Sop_i\) itself. Both must be retained when \(\G_i\) evolves.
\end{mathteach}

Subtract the two populations:
\begin{equation}
\boxed{
\mathbf K_M
=\widehat{\Qq}_{\perp}^{\top}\M^{1/2}
\left[
(\mathbf X_2-\mathbf X_1)\Vv
+\mathbf R_0\Ww(\mathbf K_{S,2}-\mathbf K_{S,1})
\right]
\boldsymbol\Gamma_0^{-1/2}.}
\label{eq:relativeCommonMetricK}
\end{equation}

\subsection{The invariant local snowball theorem}

The bases \(\widehat{\Qq}_i(t)\) now live in one ordinary Euclidean space, so the projector derivation used for \cref{eq:localSnowballTheorem} applies without alteration. Therefore:

\begin{theorem}[Common-metric local geometric snowball]
If the starting channel is separated from its complement, \(\G_i(t)\) and \(\Hh_i(t)\) are differentiable, and both populations share the same starting \(\G_0\) and \(\Hh_0\), then the common-metric chordal distance between their phenotype-space channels satisfies
\begin{equation}
\boxed{
 d_{C,M}^2(t)
=t^2\normF{\mathbf K_M}^2+O(t^3),}
\label{eq:commonMetricLocalSnowball}
\end{equation}
where \(\mathbf K_M\) is given by \cref{eq:relativeCommonMetricK}.
\end{theorem}

\begin{geombox}
There are two ways for a phenotype-space channel to turn. Its preferred eigenspace can rotate inside whitened coordinates, and the genetic ellipsoid used to unwhiten that eigenspace can itself rotate or deform. \Cref{eq:relativeCommonMetricK} adds those two tangent motions before measuring the part that leaves the ancestral channel.
\end{geombox}

\begin{biobox}
When standing allele frequencies change, the geometry of \(\G\) contributes twice: it changes the response operator \(\Sop\), and it changes how an eigenvector of \(\Sop\) is expressed as a phenotypic direction. The fully invariant susceptibility therefore distinguishes genuine evolutionary rotation from an artefact of comparing population-specific whitening systems.
\end{biobox}

\subsection{Recovery of the simpler coefficient}

If \(\G_1(t)=\G_2(t)=\G_0\) locally and choose \(\M=\G_0^{-1}\), then
\begin{equation}
\mathbf X_2-\mathbf X_1=\Zero,
\qquad
\M^{1/2}\mathbf R_0=\Ii,
\qquad
\boldsymbol\Gamma_0=\Ii_k,
\qquad
\widehat{\Qq}_{\perp}=\Ww.
\end{equation}
Consequently,
\begin{align}
\mathbf K_M
&=\Ww^\top\Ww(\mathbf K_{S,2}-\mathbf K_{S,1})\\
&=\mathbf K_{S,2}-\mathbf K_{S,1}\\
&=\cK.
\end{align}
Thus \cref{eq:commonMetricLocalSnowball} reduces exactly to \cref{eq:localSnowballTheorem}.

\begin{connectionbox}
The partition of \(\dot\Sop\) into \(\dot\G\)- and \(\dot\Hh\)-driven components, derived later in \cref{eq:SdotPartition}, determines \(\mathbf K_{S,i}\). \Cref{eq:relativeCommonMetricK} then adds the direct map-back term \((\mathbf X_2-\mathbf X_1)\Vv\). Together, those equations provide the complete first-order decomposition of channel divergence when both genetic covariance and fitness curvature evolve.
\end{connectionbox}

Because the common-metric principal angles were proven invariant to any nonsingular coordinate transformation in the preceding chapter, their derivative and the coefficient \(\normF{\mathbf K_M}^2\) are invariant as well. Coordinate changes may alter the matrices appearing in \cref{eq:relativeCommonMetricK}, but not the resulting geometric susceptibility.

\chapter{Spectral gaps, commutators, and global perturbation bounds}

\section{Why a gap alone cannot rotate a channel}

Define the commutator
\begin{equation}
[\Sop_0,\Aa]=\Sop_0\Aa-\Aa\Sop_0.
\end{equation}
If \([\Sop_0,\Aa]=\Zero\) and both matrices are symmetric, they share an orthonormal eigenbasis. In that basis, \(\Aa\) is diagonal, so
\begin{equation}
\mathbf w_a^\top\Aa\mathbf v_b=0
\end{equation}
for every cross-channel pair. Consequently
\begin{equation}
\cK(\Aa)=\Zero.
\end{equation}

\begin{geombox}
A commuting perturbation stretches or compresses the same axes but does not turn them. A noncommuting perturbation changes the orientation of the axes.
\end{geombox}

\begin{biobox}
An environment can alter the strength of selection along existing trait combinations without changing which combinations define the response channel. Geometric isolation requires selection or covariance evolution that changes trait combinations, not only response magnitudes.
\end{biobox}

\section{Davis--Kahan bound}

Let \(\widetilde\Sop=\Sop+\Delta\Sop\) be symmetric. Let \(\mathcal C\) and \(\widetilde{\mathcal C}\) be corresponding invariant subspaces separated from the rest of the spectrum by a gap \(\delta>0\). A standard Davis--Kahan sin-theta bound gives, under the usual separation convention,
\begin{equation}
\boxed{
\normtwo{\sin\boldsymbol\Theta}
\le\frac{\normtwo{\Delta\Sop}}{\delta}.}
\label{eq:DavisKahan}
\end{equation}
Here \(\normtwo{\sin\boldsymbol\Theta}=\max_j\sin\psi_j\).

\inwords The largest channel angle cannot exceed the size of the operator perturbation divided by the spectral separation protecting the channel.

\begin{cautionbox}
The Davis--Kahan expression is a bound, not the predicted angle and not the exact local coefficient. It ignores the orientation of \(\Delta\Sop\). The local formula \cref{eq:relativeK} is sharper because it retains the cross-channel matrix elements.
\end{cautionbox}

For a rank-one channel, the weakest protecting separation is
\begin{equation}
\delta_1=\lambda_1-\lambda_2.
\end{equation}
For a rank-\(k\) channel, it is
\begin{equation}
\delta_k=\lambda_k-\lambda_{k+1}.
\end{equation}

\section{A two-dimensional diagonal bound}

Suppose
\begin{equation}
\G=\diag(g_{\parallel},g_{\perp}),
\qquad
\Hh_1=\Ii,
\qquad
\Hh_2=\Ii+\diag(0,\Delta h).
\end{equation}
Then
\begin{equation}
\Delta\Sop
=\G^{1/2}(\Hh_2-\Hh_1)\G^{1/2}
=\diag(0,g_{\perp}\Delta h).
\end{equation}
Therefore
\begin{equation}
\normtwo{\Delta\Sop}=g_{\perp}|\Delta h|.
\end{equation}
If the original gap is \(\delta\),
\begin{equation}
\sin\theta\le\frac{g_{\perp}|\Delta h|}{\delta}.
\end{equation}

\inwords Variation in the direction being perturbed scales the effective operator change. A curvature change in a direction with almost no additive variance has little effect on the response channel.

\begin{cautionbox}
In this diagonal example, \(\Sop_0\) and \(\Delta\Sop\) commute. Their eigenvectors therefore do not rotate at all: the observed angle is exactly zero even though the right-hand side of the Davis--Kahan inequality can be positive. The example deliberately shows why a global upper bound is not a mechanistic prediction. The commutator and the cross-channel numerator in \cref{eq:relativeK} recover the exact zero rotation.
\end{cautionbox}

\begin{figure}[H]
\centering
\includegraphics[width=0.90\textwidth]{figures/figS2_DK_bound_checked.pdf}
\caption{The spectral gap controls the upper envelope of rotation, but observed rotation also depends on the direction of the perturbation.}
\label{fig:DavisKahanCheck}
\end{figure}

\section{Additivity and degeneracy}

Define a leading-mode concentration index
\begin{equation}
\cA_g=\frac{\lambda_1}{\sum_{j=1}^n\lambda_j}
\end{equation}
and rank-one degeneracy
\begin{equation}
\cD=\frac{\lambda_2}{\lambda_1}.
\end{equation}
In two dimensions,
\begin{align}
\cA_g
&=\frac{\lambda_1}{\lambda_1+\lambda_2}\\
&=\frac{1}{1+\cD},
\end{align}
so
\begin{equation}
\boxed{\cD=\frac{1-\cA_g}{\cA_g}.}
\end{equation}
In higher dimensions, \(1-\cA_g\) can be distributed across many modes, so low \(\cA_g\) does not identify which competitor threatens the leading channel. The explicit gap remains the cleaner stability quantity.

\begin{connectionbox}
This distinction mirrors the earlier warning about channel dimension. A scalar concentration index summarizes how dominant the first mode is overall; perturbation sensitivity depends on the particular neighbouring modes to which the perturbation couples.
\end{connectionbox}

\chapter{Deterministic selection and stochastic channel diffusion have different exponents}

\section{Deterministic divergence}

From \cref{eq:localSnowballTheorem,eq:commonMetricLocalSnowball}, smooth deterministic forcing gives
\begin{equation}
d_{C,M}^2(t)=\chi_M^2t^2+O(t^3),
\qquad
\chi_M^2=\normF{\mathbf K_M}^2.
\end{equation}
When the whitening map is shared and fixed, \(\mathbf K_M=\cK\) and \(\chi_M^2=\chi_k^2=\normF{\cK}^2\). Equivalently, each principal angle is \(O(t)\).

\section{Brownian motion in the tangent space}

Let a local \emph{normal-coordinate} tangent displacement obey
\begin{equation}
\dd\mathbf K_t=\mathbf B_K\,\dd t+\boldsymbol\Sigma_K\,\dd\mathbf W_t,
\label{eq:tangentSDE}
\end{equation}
where \(\mathbf W_t\) is a vector of independent Brownian motions after vectorizing the tangent matrix. Starting from zero,
\begin{equation}
\mathbf K_t=\mathbf B_Kt+\boldsymbol\Sigma_K\mathbf W_t.
\end{equation}

First define the quadratic tangent-space surrogate
\begin{equation}
\widetilde d_C^2(t)=\normF{\mathbf K_t}^2.
\label{eq:tangentDistanceSurrogate}
\end{equation}
Its expectation can be calculated exactly:
\begin{align}
\E[\widetilde d_C^2(t)]
&=\E\left[\normF{\mathbf B_Kt+\boldsymbol\Sigma_K\mathbf W_t}^2\right]\\
&=\normF{\mathbf B_K}^2t^2
+2t\E\innerF{\mathbf B_K}{\boldsymbol\Sigma_K\mathbf W_t}
+\E\normF{\boldsymbol\Sigma_K\mathbf W_t}^2.
\end{align}
Because \(\E[\mathbf W_t]=\Zero\), the cross term is zero. Since \(\Cov(\mathbf W_t)=t\Ii\),
\begin{equation}
\E\normF{\boldsymbol\Sigma_K\mathbf W_t}^2
=t\tr(\boldsymbol\Sigma_K\boldsymbol\Sigma_K^\top).
\end{equation}
Therefore
\begin{equation}
\boxed{
\E[\widetilde d_C^2(t)]
=a t+b t^2,}
\label{eq:selectionDriftMixture}
\end{equation}
where
\begin{equation}
a=\tr(\boldsymbol\Sigma_K\boldsymbol\Sigma_K^\top),
\qquad
b=\normF{\mathbf B_K}^2.
\end{equation}

\subsection{Relate the tangent surrogate to actual chordal distance}

Let \(\kappa_j(t)\) be the singular values of the normal-coordinate tangent matrix \(\mathbf K_t\). Along the corresponding Grassmann exponential map, the principal angles are \(\kappa_j\), so
\begin{align}
d_C^2(t)
&=\sum_j\sin^2\kappa_j(t)\\
&=\sum_j\left[\kappa_j(t)^2-\frac13\kappa_j(t)^4
+O(\kappa_j(t)^6)\right]\\
&=\normF{\mathbf K_t}^2+O(\normF{\mathbf K_t}^4).
\label{eq:stochasticChordalExpansion}
\end{align}
For Brownian motion, \(\mathbf K_t=O_p(t^{1/2})\), so the geometric correction has expectation of order \(t^2\). Consequently
\begin{equation}
\boxed{
\E[d_C^2(t)]=a t+O(t^2).}
\label{eq:stochasticLeadingLaw}
\end{equation}
The coefficient of the leading linear term is still exactly
\begin{equation}
a=\tr(\boldsymbol\Sigma_K\boldsymbol\Sigma_K^\top).
\end{equation}

When diffusion is absent, \(a=0\), the deterministic tangent term gives the quadratic coefficient \(b=\normF{\mathbf B_K}^2\), consistent with \cref{eq:localSnowballTheorem,eq:commonMetricLocalSnowball}. When drift and diffusion act together, fitting
\begin{equation}
\E[d_C^2(t)]\approx a t+b_{\mathrm{eff}}t^2
\end{equation}
is useful, but \(b_{\mathrm{eff}}\) contains not only deterministic directional forcing. It can also contain second-order manifold curvature, diffusion, and nonlinear path corrections.

\inwords Unbiased stochastic motion produces root-time angle and therefore a linear leading expectation for squared distance. Repeatable directional forcing produces linear angle and therefore a quadratic contribution. The linear coefficient cleanly diagnoses diffusion in this local model; the empirical quadratic coefficient is a composite unless stochastic second-order corrections are modelled explicitly.

\begin{mathteach}
\Cref{eq:selectionDriftMixture} is exact for the tangent-space surrogate. \Cref{eq:stochasticLeadingLaw} is the corresponding asymptotic statement for the actual curved Grassmann geometry. Keeping those statements separate prevents a second-order approximation from being mistaken for an exact decomposition of the observed chordal curve.
\end{mathteach}

\begin{biobox}
Replicate divergence is essential. Deterministic ecological forcing should produce a repeatable direction of channel rotation across comparable population pairs. Drift or mutation-order effects can generate similar magnitudes but less directional repeatability and a stronger linear baseline. Estimating the higher-order deterministic component requires either an explicit stochastic channel model or enough replication to quantify the diffusion correction.
\end{biobox}

\section{General scaling law}

If a characteristic angle obeys
\begin{equation}
\psi(t)\propto t^\alpha,
\end{equation}
then every local quadratic overlap deficit obeys
\begin{equation}
I(t)\propto t^{2\alpha}.
\end{equation}
A superlinear isolation curve requires \(\alpha>1/2\). The canonical deterministic geometric snowball has \(\alpha=1\), hence exponent two.

\begin{figure}[H]
\centering
\includegraphics[width=0.82\textwidth]{figures/02_selection_vs_drift_scaling.pdf}
\caption{Numerical comparison of deterministic and Brownian channel divergence. The fitted log--log slopes approach two and one, respectively.}
\label{fig:selectionDriftScaling}
\end{figure}

\begin{summarybox}
A quadratic curve is a prediction of smooth directional channel rotation, not a generic consequence of any divergence. A linear leading term identifies local diffusion in this model. A fitted form \(a t+b_{\mathrm{eff}}t^2\) is useful, but interpreting \(b_{\mathrm{eff}}\) as purely deterministic requires an explicit treatment of second-order stochastic and manifold corrections.
\end{summarybox}

\part{From channel geometry to reproductive isolation}

\chapter{Why angle is not yet fitness}

\section{A geometric discrepancy has no inherent biological sign}

The quantities
\begin{equation}
d_C^2=\sum_j\sin^2\psi_j
\qquad\text{and}\qquad
I_V=1-\prod_j\cos^2\psi_j
\end{equation}
measure channel discrepancy. Neither contains a fitness scale, a habitat, a hybrid mean, or a hybrid covariance. Consequently neither can, by itself, determine whether hybrids are less fit, equally fit, or more fit than their parents.

\begin{cautionbox}
Calling \(I_V\) ``reproductive isolation'' without a fitness model would be circular. The channels would be declared incompatible because they overlap poorly, and poor overlap would then be treated as evidence of incompatibility. The missing bridge is a model of hybrid phenotypes evaluated on parental fitness surfaces.
\end{cautionbox}

\section{Hybrid phenotype distribution}

Let a hybrid phenotype be a random vector
\begin{equation}
\mathbf Z_H\sim(\muv_H,\Pp_H),
\end{equation}
where
\begin{equation}
\E[\mathbf Z_H]=\muv_H,
\qquad
\Cov(\mathbf Z_H)=\Pp_H.
\end{equation}
We do not initially require normality.

In environment \(e\), let the optimum and curvature be \(\thetav_e\) and \(\Hh_e\). Define the hybrid mean displacement
\begin{equation}
\mathbf d_e=\muv_H-\thetav_e.
\end{equation}
Write
\begin{equation}
\mathbf Z_H-\thetav_e
=\mathbf d_e+\boldsymbol\varepsilon,
\end{equation}
where
\begin{equation}
\E[\boldsymbol\varepsilon]=\Zero,
\qquad
\Cov(\boldsymbol\varepsilon)=\Pp_H.
\end{equation}

\section{Three different fitness summaries}

It is useful to distinguish three quantities.

\subsection{Individual log-fitness loss}
For one phenotype,
\begin{equation}
\ell_e(\mathbf z)
=-\log\frac{W_e(\mathbf z)}{W_{0e}}
=\half(\mathbf z-\thetav_e)^\top\Hh_e(\mathbf z-\thetav_e).
\end{equation}

\subsection{Expected log-fitness loss}
\begin{equation}
L^{\log}_{H,e}=\E[\ell_e(\mathbf Z_H)].
\end{equation}
This is algebraically simple and additive across independent log-fitness components.

\subsection{Arithmetic mean fitness}
\begin{equation}
\bar W_{H,e}=\E[W_e(\mathbf Z_H)].
\end{equation}
This is often the directly relevant demographic quantity, but its relationship to expected log fitness is nonlinear.

By Jensen's inequality,
\begin{equation}
-\log\E\left[\frac{W}{W_0}\right]
\le
\E\left[-\log\frac{W}{W_0}\right],
\end{equation}
because \(-\log x\) is convex. Thus the log of arithmetic mean fitness and expected log-fitness loss are not identical.

\begin{mathteach}
The geometric mean relative fitness is \(\exp[-L^{\log}]\). The arithmetic mean relative fitness is \(\bar W/W_0\). They agree only under special conditions or approximately when variation in log fitness is small.
\end{mathteach}

\section{Bounded isolation measures}

A log-scale barrier may be mapped to \([0,1)\) by
\begin{equation}
I_L=1-e^{-\Delta L^{\log}},
\end{equation}
where \(\Delta L^{\log}\) is hybrid log-load relative to a chosen parental benchmark. This is a geometric-mean interpretation.

An arithmetic-fitness isolation measure is
\begin{equation}
I_A=1-\frac{\bar W_H}{\bar W_{\mathrm{ref}}}.
\end{equation}
For weak load, both satisfy
\begin{equation}
I_L\approx\Delta L^{\log},
\qquad
I_A\approx-\log\left(\frac{\bar W_H}{\bar W_{\mathrm{ref}}}\right),
\end{equation}
but they should not be silently interchanged.

\chapter{Expected negative log fitness: every step}

\section{Expand the hybrid quadratic form}

The individual loss is
\begin{equation}
\ell_e(\mathbf Z_H)
=\half(\mathbf d_e+\boldsymbol\varepsilon)^\top
\Hh_e
(\mathbf d_e+\boldsymbol\varepsilon).
\end{equation}
Expand the product:
\begin{align}
(\mathbf d_e+\boldsymbol\varepsilon)^\top\Hh_e(\mathbf d_e+\boldsymbol\varepsilon)
={}&\mathbf d_e^\top\Hh_e\mathbf d_e
+\mathbf d_e^\top\Hh_e\boldsymbol\varepsilon\\
&+\boldsymbol\varepsilon^\top\Hh_e\mathbf d_e
+\boldsymbol\varepsilon^\top\Hh_e\boldsymbol\varepsilon.
\end{align}
Because \(\Hh_e\) is symmetric, the middle two scalar terms are equal:
\begin{equation}
\mathbf d_e^\top\Hh_e\boldsymbol\varepsilon
=\boldsymbol\varepsilon^\top\Hh_e\mathbf d_e.
\end{equation}
Thus
\begin{equation}
\ell_e
=\half\mathbf d_e^\top\Hh_e\mathbf d_e
+\mathbf d_e^\top\Hh_e\boldsymbol\varepsilon
+\half\boldsymbol\varepsilon^\top\Hh_e\boldsymbol\varepsilon.
\end{equation}
Take expectations:
\begin{align}
L^{\log}_{H,e}
={}&\half\mathbf d_e^\top\Hh_e\mathbf d_e
+\mathbf d_e^\top\Hh_e\E[\boldsymbol\varepsilon]
+\half\E[\boldsymbol\varepsilon^\top\Hh_e\boldsymbol\varepsilon]\\
={}&\half\mathbf d_e^\top\Hh_e\mathbf d_e
+\half\tr(\Hh_e\Pp_H),
\end{align}
using \(\E[\boldsymbol\varepsilon]=\Zero\) and \cref{eq:traceQuadraticExpectation}.

Therefore
\begin{equation}
\boxed{
L^{\log}_{H,e}
=\half\mathbf d_e^\top\Hh_e\mathbf d_e
+\half\tr(\Hh_e\Pp_H).}
\label{eq:expectedLogFitnessLoad}
\end{equation}

\begin{mathteach}
The expected value of a quadratic function depends only on the first two moments of the distribution. Normality is not required for \cref{eq:expectedLogFitnessLoad}.
\end{mathteach}

\begin{geombox}
The first term evaluates the location of the hybrid cloud's centre. The second evaluates the size and orientation of the cloud itself. A cloud can be centred exactly at the optimum and still have low fitness if it extends far into steep directions.
\end{geombox}

\begin{biobox}
Hybrid failure can therefore arise from two geometrically distinct causes: mean mismatch and variance mismatch. Mean mismatch is an intermediate or transgressive phenotype lying away from the optimum. Variance mismatch is segregation or recombination placing offspring in trait combinations strongly penalised by selection.
\end{biobox}

\section{Diagonalize the curvature-weighted covariance}

Use cyclic trace invariance:
\begin{equation}
\tr(\Hh_e\Pp_H)
=\tr(\Hh_e^{1/2}\Pp_H\Hh_e^{1/2}).
\end{equation}
Let
\begin{equation}
\mathbf C_e=\Hh_e^{1/2}\Pp_H\Hh_e^{1/2}.
\end{equation}
Since \(\mathbf C_e\) is symmetric positive semidefinite, let its eigenvalues be \(c_{ej}\ge0\). Then
\begin{equation}
\tr(\Hh_e\Pp_H)=\sum_{j=1}^nc_{ej}.
\end{equation}

\inwords Each \(c_{ej}\) is hybrid variance measured after scaling phenotype directions by the square root of fitness curvature. The trace is the total curvature-standardised hybrid variance.

\begin{connectionbox}
This has exactly the same algebraic structure as the channel operator. The response operator \(\G^{1/2}\Hh\G^{1/2}\) asks how curvature acts on parental additive variation. The hybrid-load matrix \(\Hh^{1/2}\Pp_H\Hh^{1/2}\) asks how hybrid phenotypic variance is penalised by curvature. The order of the square-root factors reflects the question being asked, but the nonzero eigenvalues of \(\G^{1/2}\Hh\G^{1/2}\) and \(\Hh^{1/2}\G\Hh^{1/2}\) coincide.
\end{connectionbox}

\chapter{Exact arithmetic mean fitness under Gaussian phenotypes}

\section{Set up the Gaussian integral}

Now assume
\begin{equation}
\mathbf Z_H\sim N(\muv_H,\Pp_H).
\end{equation}
Define
\begin{equation}
\mathbf Y=\Hh_e^{1/2}(\mathbf Z_H-\thetav_e).
\end{equation}
Then
\begin{equation}
\mathbf Y\sim N(\mathbf m,\mathbf C),
\end{equation}
where
\begin{equation}
\mathbf m=\Hh_e^{1/2}\mathbf d_e,
\qquad
\mathbf C=\Hh_e^{1/2}\Pp_H\Hh_e^{1/2}.
\end{equation}
Relative fitness is
\begin{equation}
\frac{W}{W_0}=e^{-\half\mathbf Y^\top\mathbf Y}.
\end{equation}
Therefore
\begin{equation}
\frac{\bar W_{H,e}}{W_{0e}}
=\int_{\R^n}e^{-\half\mathbf y^\top\mathbf y}
\frac{e^{-\half(\mathbf y-\mathbf m)^\top\mathbf C^{-1}(\mathbf y-\mathbf m)}}
{(2\pi)^{n/2}|\mathbf C|^{1/2}}\,\dd\mathbf y.
\label{eq:GaussianFitnessIntegral}
\end{equation}
For clarity, first assume \(\mathbf C\succ0\). The semidefinite case follows by continuity or restriction to its support.

\section{Combine the exponents}

Expand the Gaussian exponent:
\begin{align}
(\mathbf y-\mathbf m)^\top\mathbf C^{-1}(\mathbf y-\mathbf m)
={}&\mathbf y^\top\mathbf C^{-1}\mathbf y
-2\mathbf m^\top\mathbf C^{-1}\mathbf y
+\mathbf m^\top\mathbf C^{-1}\mathbf m.
\end{align}
Adding \(\mathbf y^\top\mathbf y\), the total exponent is
\begin{align}
&-\half\left[
\mathbf y^\top(\mathbf C^{-1}+\Ii)\mathbf y
-2\mathbf m^\top\mathbf C^{-1}\mathbf y
+\mathbf m^\top\mathbf C^{-1}\mathbf m
\right].
\end{align}
Define
\begin{equation}
\mathbf A=\mathbf C^{-1}+\Ii,
\qquad
\mathbf b=\mathbf C^{-1}\mathbf m.
\end{equation}
Then
\begin{equation}
\mathbf y^\top\mathbf A\mathbf y-2\mathbf b^\top\mathbf y
=(\mathbf y-\mathbf A^{-1}\mathbf b)^\top\mathbf A
(\mathbf y-\mathbf A^{-1}\mathbf b)
-\mathbf b^\top\mathbf A^{-1}\mathbf b.
\label{eq:completeSquareMatrix}
\end{equation}

\why Expanding the right-hand side gives \(\mathbf y^\top\mathbf A\mathbf y-2\mathbf b^\top\mathbf y+\mathbf b^\top\mathbf A^{-1}\mathbf b-\mathbf b^\top\mathbf A^{-1}\mathbf b\), recovering the left-hand side.

Substitute into the integral. Terms independent of \(\mathbf y\) can be moved outside:
\begin{align}
\frac{\bar W_{H,e}}{W_{0e}}
={}&|\mathbf C|^{-1/2}
\exp\left[-\half\left(
\mathbf m^\top\mathbf C^{-1}\mathbf m
-\mathbf b^\top\mathbf A^{-1}\mathbf b
\right)\right]\\
&\times
(2\pi)^{-n/2}
\int
\exp\left[-\half(\mathbf y-\mathbf A^{-1}\mathbf b)^\top
\mathbf A
(\mathbf y-\mathbf A^{-1}\mathbf b)
\right]\dd\mathbf y.
\end{align}
The remaining Gaussian integral is
\begin{equation}
(2\pi)^{n/2}|\mathbf A|^{-1/2}.
\end{equation}
Therefore
\begin{equation}
\frac{\bar W_{H,e}}{W_{0e}}
=|\mathbf C|^{-1/2}|\mathbf A|^{-1/2}
\exp\left[-\half\left(
\mathbf m^\top\mathbf C^{-1}\mathbf m
-\mathbf b^\top\mathbf A^{-1}\mathbf b
\right)\right].
\label{eq:GaussianFitnessIntermediate}
\end{equation}

\section{Simplify the determinant}

Since
\begin{equation}
\mathbf A=\mathbf C^{-1}(\Ii+\mathbf C),
\end{equation}
we have
\begin{align}
|\mathbf C||\mathbf A|
&=|\mathbf C||\mathbf C^{-1}||\Ii+\mathbf C|\\
&=|\Ii+\mathbf C|.
\end{align}
Thus
\begin{equation}
|\mathbf C|^{-1/2}|\mathbf A|^{-1/2}
=|\Ii+\mathbf C|^{-1/2}.
\end{equation}

\section{Simplify the exponent}

Because \(\mathbf b=\mathbf C^{-1}\mathbf m\),
\begin{equation}
\mathbf b^\top\mathbf A^{-1}\mathbf b
=\mathbf m^\top\mathbf C^{-1}\mathbf A^{-1}\mathbf C^{-1}\mathbf m.
\end{equation}
Also
\begin{equation}
\mathbf A^{-1}=(\Ii+\mathbf C)^{-1}\mathbf C,
\end{equation}
so
\begin{align}
\mathbf C^{-1}\mathbf A^{-1}\mathbf C^{-1}
&=\mathbf C^{-1}(\Ii+\mathbf C)^{-1}.
\end{align}
Therefore the exponent matrix is
\begin{align}
\mathbf C^{-1}-\mathbf C^{-1}(\Ii+\mathbf C)^{-1}
&=\mathbf C^{-1}\left[\Ii-(\Ii+\mathbf C)^{-1}\right]\\
&=\mathbf C^{-1}\left[\mathbf C(\Ii+\mathbf C)^{-1}\right]\\
&=(\Ii+\mathbf C)^{-1}.
\end{align}
The bracket identity follows because
\begin{equation}
\Ii-(\Ii+\mathbf C)^{-1}
=\left[(\Ii+\mathbf C)-\Ii\right](\Ii+\mathbf C)^{-1}
=\mathbf C(\Ii+\mathbf C)^{-1}.
\end{equation}

Combining determinant and exponent gives
\begin{equation}
\boxed{
\frac{\bar W_{H,e}}{W_{0e}}
=|\Ii+\mathbf C|^{-1/2}
\exp\left[-\half\mathbf m^\top(\Ii+\mathbf C)^{-1}\mathbf m\right].}
\label{eq:GaussianFitnessCompact}
\end{equation}
Substituting \(\mathbf C=\Hh_e^{1/2}\Pp_H\Hh_e^{1/2}\) and \(\mathbf m=\Hh_e^{1/2}\mathbf d_e\),
\begin{equation}
\boxed{
\begin{aligned}
\bar W_{H,e}
={}&W_{0e}
\left|\Ii+\Hh_e^{1/2}\Pp_H\Hh_e^{1/2}\right|^{-1/2}\\
&\times
\exp\left[
-\half\mathbf d_e^\top\Hh_e^{1/2}
\left(\Ii+\Hh_e^{1/2}\Pp_H\Hh_e^{1/2}\right)^{-1}
\Hh_e^{1/2}\mathbf d_e
\right].
\end{aligned}}
\label{eq:exactGaussianMeanFitness}
\end{equation}

\begin{geombox}
The determinant term is the loss caused by spreading a Gaussian cloud through curved fitness space. The exponential term is the loss caused by displacing its centre. Variance also softens the mean penalty through the inverse \((\Ii+\mathbf C)^{-1}\), because a broad distribution samples phenotypes on both sides of its centre.
\end{geombox}

\begin{cautionbox}
\Cref{eq:exactGaussianMeanFitness} gives arithmetic mean fitness. \Cref{eq:expectedLogFitnessLoad} gives expected log-fitness loss. Both are exact under their assumptions, but they answer different averaging questions.
\end{cautionbox}

\chapter{The general covariance--curvature trace identity}

\section{Define a parental benchmark}

Assume parental population \(i\) is centred at its own optimum and has covariance \(\G_i\). Its covariance component of expected log loss is
\begin{equation}
L_{P_i,i}^{\mathrm{cov}}=\half\tr(\Hh_i\G_i).
\end{equation}
The hybrid covariance loss in environment \(i\) is
\begin{equation}
L_{H,i}^{\mathrm{cov}}=\half\tr(\Hh_i\Pp_H).
\end{equation}
The hybrid excess relative to that parent is
\begin{equation}
\Delta L_i^{\mathrm{cov}}
=\half\tr[\Hh_i(\Pp_H-\G_i)].
\end{equation}
Take an equal average over the two parental environments:
\begin{equation}
L_{\mathrm{cov}}
=\half(\Delta L_1^{\mathrm{cov}}+\Delta L_2^{\mathrm{cov}})
=\frac14\left
\{\tr[\Hh_1(\Pp_H-\G_1)]
+\tr[\Hh_2(\Pp_H-\G_2)]\right\}.
\label{eq:covLoadBenchmark}
\end{equation}

\section{Decompose hybrid covariance}

Write
\begin{equation}
\boxed{\Pp_H=\half(\G_1+\G_2)+\Rr.}
\label{eq:hybridCovDecomposition}
\end{equation}
The matrix \(\Rr\) is excess covariance beyond the parental average. It can represent segregation variance, linkage disequilibrium, transgressive covariance, environmental variance differences, or any other departure from the simple average, depending on the biological model.

Then
\begin{align}
\Pp_H-\G_1
&=\half(\G_2-\G_1)+\Rr,\\
\Pp_H-\G_2
&=-\half(\G_2-\G_1)+\Rr.
\end{align}
Substitute into \cref{eq:covLoadBenchmark}:
\begin{align}
L_{\mathrm{cov}}
={}&\frac14\left\{
\half\tr[\Hh_1(\G_2-\G_1)]
+\tr(\Hh_1\Rr)\right.\\
&\left.\qquad
-\half\tr[\Hh_2(\G_2-\G_1)]
+\tr(\Hh_2\Rr)
\right\}\\
={}&\frac18\tr[(\Hh_1-\Hh_2)(\G_2-\G_1)]
+\frac14\tr[(\Hh_1+\Hh_2)\Rr]\\
={}&-\frac18\tr[(\Hh_2-\Hh_1)(\G_2-\G_1)]
+\frac14\tr[(\Hh_1+\Hh_2)\Rr].
\end{align}
Therefore
\begin{theorem}[General covariance--curvature identity]
Under \cref{eq:hybridCovDecomposition}, the equally weighted hybrid covariance log-load relative to the parents in their own environments is
\begin{equation}
\boxed{
L_{\mathrm{cov}}
=-\frac18\tr[(\Hh_2-\Hh_1)(\G_2-\G_1)]
+\frac14\tr[(\Hh_1+\Hh_2)\Rr].}
\label{eq:generalTraceIdentity}
\end{equation}
\end{theorem}

\section{Interpret the sign}

The first term is a negative Frobenius inner product:
\begin{equation}
-\frac18\innerF{\Hh_2-\Hh_1}{\G_2-\G_1}.
\end{equation}
It is positive when directions that gain genetic variance in population 2 tend to lose curvature there relative to population 1, and vice versa. It is negative under the opposite association.

\begin{biobox}
Positive geometric hybrid load requires covariance--curvature coadaptation: each population carries relatively more variance in directions that are relatively less costly in its own environment. Recombining those covariance structures then places hybrid variance into directions that are not jointly protected by either fitness surface.
\end{biobox}

\begin{cautionbox}
Large channel angle alone does not determine the sign of the trace. If variance and curvature differences are aligned in the opposite sense, the covariance term can be negative and hybrids can exhibit heterosis relative to the chosen parental benchmark.
\end{cautionbox}

\section{Interpret segregation excess}

If \(\Rr\succeq0\), then
\begin{equation}
\tr[(\Hh_1+\Hh_2)\Rr]\ge0
\end{equation}
because both matrices are positive semidefinite. Thus additional segregation variance raises expected log load. If \(\Rr\) is indefinite, some recombination-generated covariances can instead reduce load.

\inwords Segregation excess is not automatically an isotropic inflation. Its orientation matters. Variance added along steep directions is costly; variance added along shallow directions can be comparatively benign.

\chapter{The channel-ridge theorem: exact calibration by principal angles}

\section{Structured parental covariance and curvature}

The theorem is written in one fixed reference metric. Equivalently, first transform phenotype space by \(\M^{1/2}\) so that the chosen common metric becomes Euclidean, and then define ordinary orthogonal projectors in that transformed space. Let \(\Pp_i\) be the resulting rank-\(k\) projector onto population \(i\)'s channel. Define
\begin{align}
\G_i
&=g_{\perp}\Ii+(g_{\parallel}-g_{\perp})\Pp_i,
\label{eq:ridgeG}\\
\Hh_i
&=h_{\perp}\Ii-(h_{\perp}-h_{\parallel})\Pp_i.
\label{eq:ridgeH}
\end{align}
Set
\begin{equation}
\Delta g=g_{\parallel}-g_{\perp}>0,
\qquad
\Delta h=h_{\perp}-h_{\parallel}>0.
\end{equation}
Then
\begin{equation}
\G_i=g_{\perp}\Ii+\Delta g\Pp_i,
\qquad
\Hh_i=h_{\perp}\Ii-\Delta h\Pp_i.
\end{equation}

\begin{cautionbox}
The result is invariant to a change of trait units when the metric is transformed with those units, but it can depend on the scientifically chosen reference metric itself. In applications with different \(\G_i\), the projectors should therefore be constructed from the common-metric bases of \cref{eq:metricEuclideanization,eq:EuclideanizedChannelBasis}, not from raw population-specific whitened eigenvectors.
\end{cautionbox}

\inwords Genetic variance is enriched along the channel, while fitness curvature is reduced along the same channel. Each parental covariance is therefore coadapted to a shallow local ridge.

For \(\Pp_i\) to be the leading response subspace of \(\Sop_i\), require
\begin{equation}
g_{\parallel}h_{\parallel}>g_{\perp}h_{\perp}.
\label{eq:ridgeLeadingCondition}
\end{equation}
Indeed, along \(\Pp_i\) the response eigenvalue is \(g_{\parallel}h_{\parallel}\); in the complement it is \(g_{\perp}h_{\perp}\).

\section{Compute covariance and curvature differences}

Subtract the two covariance matrices:
\begin{equation}
\G_2-\G_1=\Delta g(\Pp_2-\Pp_1).
\end{equation}
Subtract the curvature matrices:
\begin{equation}
\Hh_2-\Hh_1=-\Delta h(\Pp_2-\Pp_1).
\end{equation}
With \(\Rr=\Zero\), substitute into \cref{eq:generalTraceIdentity}:
\begin{align}
L_{\mathrm{cov}}
&=-\frac18\tr\left[
-\Delta h(\Pp_2-\Pp_1)
\Delta g(\Pp_2-\Pp_1)
\right]\\
&=\frac{\Delta g\Delta h}{8}
\tr[(\Pp_2-\Pp_1)^2].
\end{align}
Use \cref{eq:chordalIdentity}:
\begin{equation}
\tr[(\Pp_2-\Pp_1)^2]
=\normF{\Pp_2-\Pp_1}^2
=2\sum_{j=1}^k\sin^2\psi_j.
\end{equation}
Therefore
\begin{theorem}[Channel-ridge fitness calibration]
Under \cref{eq:ridgeG,eq:ridgeH}, equal channel rank, zero segregation excess, and the parental benchmark above,
\begin{equation}
\boxed{
L_{\mathrm{cov}}
=\frac{\Delta g\Delta h}{4}
\sum_{j=1}^k\sin^2\psi_j
=\frac{\Delta g\Delta h}{4}d_C^2.}
\label{eq:ridgeTheorem}
\end{equation}
\end{theorem}

\inwords Channel mismatch becomes a positive fitness load because the exact same projector difference controls both covariance divergence and curvature divergence, with opposite signs.

\section{Early-time fitness snowball}

Combine \cref{eq:ridgeTheorem} with \cref{eq:localSnowballTheorem}:
\begin{align}
L_{\mathrm{cov}}(t)
&=\frac{\Delta g\Delta h}{4}
\left[t^2\normF{\cK}^2+O(t^3)\right]\\
&=\boxed{
\frac{\Delta g\Delta h}{4}\normF{\cK}^2t^2+O(t^3).}
\label{eq:fitnessSnowballLocal}
\end{align}

For equal principal angles \(\psi_j=\theta\),
\begin{equation}
L_{\mathrm{cov}}
=\frac{\Delta g\Delta h}{4}k\sin^2\theta
\approx\frac{\Delta g\Delta h}{4}k\theta^2.
\end{equation}

\begin{connectionbox}
The theory now has two distinct quadratic statements. Channel discrepancy is quadratic in time under deterministic forcing. Hybrid covariance load is proportional to channel discrepancy under covariance--curvature coadaptation. Only their combination justifies calling the result a reproductive-isolation snowball.
\end{connectionbox}

\begin{figure}[H]
\centering
\includegraphics[width=\textwidth]{figures/pedagogy_hybrid_load.pdf}
\caption{Geometric reading of the channel-ridge theorem. Parental genetic variance lies along differently oriented shallow ridges. The averaged hybrid covariance occupies directions that are steep in the parental fitness geometries.}
\label{fig:hybridLoadPedagogy}
\end{figure}

\begin{figure}[H]
\centering
\includegraphics[width=0.82\textwidth]{figures/03_hybrid_load_calibration.pdf}
\caption{Numerical calibration of raw channel geometry against the exact covariance--curvature load. Geometry acquires a fitness scale only after \(\Delta g\) and \(\Delta h\) are specified.}
\label{fig:hybridLoadCheck}
\end{figure}

\chapter{Mean mismatch, segregation excess, and hybrid generation}

\section{Midpoint mean mismatch}

Suppose parental means lie at their own optima \(\thetav_1\) and \(\thetav_2\), and the hybrid mean is the midpoint
\begin{equation}
\muv_H=\half(\thetav_1+\thetav_2).
\end{equation}
Define
\begin{equation}
\Delta\thetav=\thetav_2-\thetav_1.
\end{equation}
Then in environment 1,
\begin{equation}
\mathbf d_1=\muv_H-\thetav_1=\half\Delta\thetav.
\end{equation}
In environment 2,
\begin{equation}
\mathbf d_2=\muv_H-\thetav_2=-\half\Delta\thetav.
\end{equation}
The mean log-loads are
\begin{align}
L_{\mathrm{mean},1}
&=\half\mathbf d_1^\top\Hh_1\mathbf d_1
=\frac18\Delta\thetav^\top\Hh_1\Delta\thetav,\\
L_{\mathrm{mean},2}
&=\frac18\Delta\thetav^\top\Hh_2\Delta\thetav.
\end{align}
Their equal average is
\begin{equation}
\boxed{
L_{\mathrm{mean}}
=\frac1{16}\Delta\thetav^\top(\Hh_1+\Hh_2)\Delta\thetav.}
\label{eq:midpointMeanLoad}
\end{equation}

\inwords Mean mismatch depends on phenotypic displacement between optima and the curvature along that displacement. It is not a channel-angle term and should be estimated separately.

\section{A complete log-load decomposition}

For the equal parental-environment benchmark, write
\begin{equation}
L_E=L_{\mathrm{mean}}+L_{\mathrm{cov}}.
\end{equation}
Using \cref{eq:generalTraceIdentity},
\begin{align}
L_E
={}&L_{\mathrm{mean}}
-\frac18\tr[(\Hh_2-\Hh_1)(\G_2-\G_1)]\\
&+\frac14\tr[(\Hh_1+\Hh_2)\Rr].
\label{eq:completeHybridLogLoad}
\end{align}
This separates:
\begin{enumerate}[leftmargin=1.6em]
\item displacement of the hybrid mean;
\item covariance--curvature coadaptation broken by hybridization;
\item segregation or other excess covariance.
\end{enumerate}

\section{Why F1 and F2 hybrids can differ sharply}

An F1 may have an approximately intermediate mean but comparatively little segregation variance because homologous parental haplotypes remain intact. Dominance can create heterosis and mask recessive incompatibilities. In F2 or later recombinant generations, recombination and segregation can increase \(\Rr\), create transgressive combinations, and expose epistatic incompatibilities.

\begin{biobox}
The additive covariance mechanism is therefore expected to be especially visible in recombinant generations. F1 fitness is not an unambiguous test: strong heterosis can coexist with a geometric load that emerges only after segregation.
\end{biobox}

\begin{connectionbox}
This generation-specific prediction matches the existing \textit{Senecio} pattern that F1 hybrids can show heterosis while F2 hybrids show strong fertility or fitness reduction, and later generations can recover after incompatible variation has been removed. The theory supplies a quantitative route for separating mean, covariance, dominance, and intrinsic interaction components rather than treating all hybrid outcomes as one phenomenon.
\end{connectionbox}

\section{Environment dependence distinguishes early ecological load from intrinsic load}

The terms containing \(\Hh_e\) depend on the environment in which fitness is measured. If hybrids fail only in parental habitats, the barrier is extrinsic. If they fail in a benign common environment, an intrinsic component is implicated.

\begin{summarybox}
A decisive experiment assays the same recombinant families in both parental habitats and in a common environment. The habitat-dependent component tests geometric ecological load; the common-environment residual tests intrinsic incompatibility.
\end{summarybox}

\part{Evolution of the channels and the handoff to intrinsic isolation}

\chapter{A multilocus route from standing variation to evolving \(\G\)}

\section{Additive allelic effects}

Consider \(L\) diploid biallelic loci. At locus \(\ell\), let the count of the focal allele be
\begin{equation}
X_\ell\in\{0,1,2\}.
\end{equation}
Let \(\mathbf a_\ell\in\R^n\) be the additive phenotypic effect per copy. Relative to a reference mean, the additive phenotype is
\begin{equation}
\mathbf z_A=\sum_{\ell=1}^L(X_\ell-2p_{\ell0})\mathbf a_\ell.
\label{eq:additiveGenotypeModel}
\end{equation}

If the allele frequency in population \(i\) is \(p_{i\ell}\), then under Hardy--Weinberg sampling
\begin{equation}
\E[X_\ell]=2p_{i\ell}.
\end{equation}
Therefore the mean additive phenotype is
\begin{align}
\zbar_i
&=\E[\mathbf z_A]\\
&=\sum_\ell(2p_{i\ell}-2p_{\ell0})\mathbf a_\ell\\
&=\boxed{2\sum_{\ell=1}^L(p_{i\ell}-p_{\ell0})\mathbf a_\ell.}
\label{eq:meanFromAlleleFrequencies}
\end{align}

\inwords The population mean is a weighted sum of allelic effect vectors. Changing an allele frequency moves the mean in the direction of that allele's multivariate effect.

\section{Derive \(\G\) under linkage equilibrium}

At a biallelic diploid locus under Hardy--Weinberg equilibrium,
\begin{equation}
\Var(X_\ell)=2p_{i\ell}(1-p_{i\ell}).
\end{equation}
The covariance contributed by that locus is
\begin{align}
\Cov(X_\ell\mathbf a_\ell)
&=\Var(X_\ell)\mathbf a_\ell\mathbf a_\ell^\top\\
&=2p_{i\ell}(1-p_{i\ell})\mathbf a_\ell\mathbf a_\ell^\top.
\end{align}
If loci are in linkage equilibrium, cross-locus covariances vanish, so contributions add:
\begin{equation}
\boxed{
\G_i=2\sum_{\ell=1}^Lp_{i\ell}(1-p_{i\ell})
\mathbf a_\ell\mathbf a_\ell^\top.}
\label{eq:GfromAlleleFrequencies}
\end{equation}

\begin{geombox}
Each outer product \(\mathbf a_\ell\mathbf a_\ell^\top\) is a rank-one covariance ellipsoid pointing along the allelic effect vector. The full \(\G\) is the weighted sum of those rank-one contributions. Alleles near frequency one half receive the largest weight.
\end{geombox}

With linkage disequilibrium, additional cross-locus terms appear. Schematically,
\begin{equation}
\G_i
=2\sum_\ell p_{i\ell}q_{i\ell}\mathbf a_\ell\mathbf a_\ell^\top
+\sum_{\ell\neq m}\Cov(X_\ell,X_m)\mathbf a_\ell\mathbf a_m^\top.
\end{equation}
The linkage-equilibrium model is therefore a transparent baseline rather than a claim that LD is biologically absent.

\section{Weak-selection allele-frequency dynamics}

Let the local directional selection gradient be \(\betav_i\). The marginal directional effect of one focal allele at locus \(\ell\) is
\begin{equation}
s_{i\ell}=\mathbf a_\ell^\top\betav_i.
\end{equation}
Under weak additive selection, the leading allele-frequency change is
\begin{equation}
\dot p_{i\ell}^{\mathrm{sel}}
=p_{i\ell}(1-p_{i\ell})\mathbf a_\ell^\top\betav_i.
\label{eq:alleleSelectionDynamics}
\end{equation}
With symmetric migration from population \(j\) at rate \(m\),
\begin{equation}
\boxed{
\dot p_{i\ell}
=p_{i\ell}(1-p_{i\ell})\mathbf a_\ell^\top\betav_i
+m(p_{j\ell}-p_{i\ell}).}
\label{eq:alleleDynamicsMigration}
\end{equation}

\section{Recover the multivariate response equation, including migration}

Differentiate \cref{eq:meanFromAlleleFrequencies}:
\begin{equation}
\dot{\zbar}_i=2\sum_\ell\dot p_{i\ell}\mathbf a_\ell.
\end{equation}
Substitute \cref{eq:alleleDynamicsMigration}:
\begin{align}
\dot{\zbar}_i
={}&2\sum_\ell p_{i\ell}(1-p_{i\ell})
\mathbf a_\ell\mathbf a_\ell^\top\betav_i\\
&+2m\sum_\ell(p_{j\ell}-p_{i\ell})\mathbf a_\ell.
\end{align}
The first sum is \(\G_i\betav_i\) by \cref{eq:GfromAlleleFrequencies}. The second is
\begin{align}
2m\sum_\ell(p_{j\ell}-p_{i\ell})\mathbf a_\ell
&=m(\zbar_j-\zbar_i),
\end{align}
provided the populations share the same effect-vector representation and reference. Thus
\begin{equation}
\boxed{
\dot{\zbar}_i
=\G_i\betav_i+m(\zbar_j-\zbar_i).}
\label{eq:LandeWithMigration}
\end{equation}

\begin{cautionbox}
When migration is included at the allele-frequency level, it produces an explicit mean-migration term. The simpler identity \(\dot{\zbar}=\G\betav\) is the selection component or the \(m=0\) limit, not the full dynamics.
\end{cautionbox}

\section{How allele-frequency change alters \(\G\)}

Differentiate \cref{eq:GfromAlleleFrequencies}:
\begin{align}
\dot\G_i
&=2\sum_\ell\frac{\dd}{\dd t}
[p_{i\ell}(1-p_{i\ell})]
\mathbf a_\ell\mathbf a_\ell^\top\\
&=2\sum_\ell[(1-p_{i\ell})-p_{i\ell}]\dot p_{i\ell}
\mathbf a_\ell\mathbf a_\ell^\top\\
&=\boxed{
2\sum_\ell(1-2p_{i\ell})\dot p_{i\ell}
\mathbf a_\ell\mathbf a_\ell^\top.}
\label{eq:GdotAlleles}
\end{align}

\inwords An allele rising toward frequency one half increases its contribution to additive variance; an allele moving from one half toward fixation decreases it. Because different loci point in different multivariate directions, unequal frequency changes can rotate \(\G\) as well as change its total magnitude.

\begin{biobox}
The mechanism does not require new mutation or drift. Divergent selection can reweight pre-existing effect vectors in different populations, causing their covariance ellipsoids and response channels to rotate apart using standing variation alone.
\end{biobox}

\begin{figure}[H]
\centering
\includegraphics[width=0.84\textwidth]{figures/05_standing_variation_rotation.pdf}
\caption{A deterministic multilocus simulation in which divergent selection changes standing allele frequencies, rotates \(\G\), and generates channel misalignment.}
\label{fig:standingVariationRotation}
\end{figure}

\chapter{Differentiating \(\G^{1/2}\) and partitioning channel rotation}

\section{Differentiate the matrix square root}

Let
\begin{equation}
\Rr=\G^{1/2},
\end{equation}
so
\begin{equation}
\Rr^2=\G.
\end{equation}
Differentiate the product:
\begin{equation}
\dot\Rr\Rr+\Rr\dot\Rr=\dot\G.
\end{equation}
Define
\begin{equation}
\mathbf X=\dot\Rr=D(\G^{1/2})[\dot\G].
\end{equation}
Then \(\mathbf X\) solves the Sylvester equation
\begin{equation}
\boxed{\Rr\mathbf X+\mathbf X\Rr=\dot\G.}
\label{eq:SylvesterSqrt}
\end{equation}

\usual A Sylvester equation has the form \(\mathbf A\mathbf X+\mathbf X\mathbf B=\mathbf C\). Here it appears because differentiating a noncommuting square \(\Rr\Rr\) produces one derivative on each side.

\section{Explicit solution in the eigenbasis of \(\G\)}

Let
\begin{equation}
\G=\mathbf U\diag(g_1,\ldots,g_n)\mathbf U^\top,
\end{equation}
so
\begin{equation}
\Rr=\mathbf U\diag(r_1,\ldots,r_n)\mathbf U^\top,
\qquad r_i=\sqrt{g_i}.
\end{equation}
Write
\begin{equation}
\mathbf X=\mathbf U\mathbf Y\mathbf U^\top,
\qquad
\mathbf B=\mathbf U^\top\dot\G\mathbf U.
\end{equation}
Substitute into \cref{eq:SylvesterSqrt} and multiply by \(\mathbf U^\top\) and \(\mathbf U\):
\begin{equation}
\diag(r_i)\mathbf Y+\mathbf Y\diag(r_i)=\mathbf B.
\end{equation}
Element \((i,j)\) satisfies
\begin{equation}
r_iY_{ij}+Y_{ij}r_j=B_{ij}.
\end{equation}
Therefore
\begin{equation}
\boxed{
Y_{ij}=\frac{B_{ij}}{r_i+r_j}
=\frac{(\mathbf U^\top\dot\G\mathbf U)_{ij}}
{\sqrt{g_i}+\sqrt{g_j}}.}
\label{eq:SqrtDerivativeElements}
\end{equation}
Finally,
\begin{equation}
\mathbf X=\mathbf U\mathbf Y\mathbf U^\top.
\end{equation}

\inwords Changes in \(\G\) are filtered by the existing square-root scale. Off-diagonal changes in the eigenbasis of \(\G\) contribute directly to rotation of the square-root deformation.

\section{Differentiate the channel operator}

Recall
\begin{equation}
\Sop=\Rr\Hh\Rr.
\end{equation}
Apply the product rule to all three factors:
\begin{align}
\dot\Sop
&=\dot\Rr\Hh\Rr
+\Rr\dot\Hh\Rr
+\Rr\Hh\dot\Rr\\
&=\boxed{
\mathbf X\Hh\Rr
+\Rr\dot\Hh\Rr
+\Rr\Hh\mathbf X.}
\label{eq:SdotPartition}
\end{align}
Define
\begin{align}
\dot\Sop_G
&=\mathbf X\Hh\Rr+\Rr\Hh\mathbf X,\\
\dot\Sop_H
&=\Rr\dot\Hh\Rr.
\end{align}
Then
\begin{equation}
\dot\Sop=\dot\Sop_G+\dot\Sop_H.
\end{equation}

\begin{biobox}
The first component is channel change caused by evolution of standing genetic variance. The second is channel change caused by evolution or environmental change in fitness curvature. Their effects can reinforce, cancel, or rotate the channel in different directions.
\end{biobox}

\section{Partition the response-operator eigenspace rotation}

The tangent map in \cref{eq:tangentMap} is linear:
\begin{equation}
\cT(\Aa+\Bb)=\cT(\Aa)+\cT(\Bb).
\end{equation}
Therefore the relative rotation of the leading eigenspace of \(\Sop\) partitions as
\begin{equation}
\Delta\mathbf K_S
=\underbrace{\cT(\Delta\dot\Sop_G)}_{\mathbf K_{S,G}}
+\underbrace{\cT(\Delta\dot\Sop_H)}_{\mathbf K_{S,H}},
\label{eq:KpartitionGH}
\end{equation}
where \(\Delta\) denotes population 2 minus population 1.

Its squared magnitude is
\begin{align}
\normF{\Delta\mathbf K_S}^2
={}&\normF{\mathbf K_{S,G}}^2
+\normF{\mathbf K_{S,H}}^2
+2\innerF{\mathbf K_{S,G}}{\mathbf K_{S,H}}.
\label{eq:susceptibilityPartition}
\end{align}

\inwords The cross term measures interaction between genetic-covariance evolution and curvature evolution inside the symmetric response operator. Positive alignment accelerates eigenspace rotation; negative alignment buffers it.

For evolving \(\G\), this is not yet the complete phenotype-space susceptibility. Substitute \(\Delta\mathbf K_S\) and the Sylvester solutions \(\mathbf X_i\) into \cref{eq:relativeCommonMetricK}:
\begin{equation}
\mathbf K_M
=\widehat{\Qq}_{\perp}^{\top}\M^{1/2}
\left[(\mathbf X_2-\mathbf X_1)\Vv
+\mathbf R_0\Ww\Delta\mathbf K_S\right]
\boldsymbol\Gamma_0^{-1/2}.
\label{eq:fullEmpiricalKM}
\end{equation}
The first term is the direct deformation of the phenotype map by \(\dot\G\); the second carries the \(\G\)- and \(\Hh\)-driven rotation of the \(\Sop\) eigenspace.

\begin{connectionbox}
This gives the complete computational bridge to empirical work. Posterior draws of \(\dot\G\) and \(\dot\Hh\) are propagated through the Sylvester equation and \(\dot\Sop\), then combined by \cref{eq:fullEmpiricalKM}. The invariant local susceptibility is \(\normF{\mathbf K_M}^2\), while \(\normF{\Delta\mathbf K_S}^2\) is the informative but partial susceptibility of the response-operator eigenspace.
\end{connectionbox}

\chapter{From an early geometric barrier to Orr's DMI snowball}

\section{Effective migration as a barrier-weighted flux}

Let \(m_0\) be the potential migration rate before selection. If an early ecological or geometric log barrier is \(L_E(t)\), a minimal transmission closure is
\begin{equation}
\boxed{
m_{\mathrm{eff}}(t)=m_0e^{-L_E(t)}.}
\label{eq:effectiveMigration}
\end{equation}
If
\begin{equation}
I_E=1-e^{-L_E},
\end{equation}
then
\begin{equation}
m_{\mathrm{eff}}=m_0(1-I_E).
\end{equation}

\inwords Only the fraction of migrant or hybrid ancestry surviving the early barrier contributes effectively to homogenizing gene flow.

\begin{cautionbox}
\Cref{eq:effectiveMigration} is a modelling closure. Its exact demographic form depends on life history, mating, density regulation, and the stage at which selection acts. Its virtue is transparency: an additive log barrier becomes multiplicative survival.
\end{cautionbox}

\section{A simple establishment function}

Let candidate lineage-specific substitutions arise at rate \(\nu_i\). Use
\begin{equation}
\pi_i(m)=\frac{s_i}{s_i+2m}
\label{eq:establishmentClosure}
\end{equation}
as a dimensionless establishment function that declines with homogenizing migration. Then
\begin{equation}
\boxed{
\dot K_i=\nu_i\pi_i(m_{\mathrm{eff}}),}
\label{eq:substitutionAccumulation}
\end{equation}
where \(K_i\) is the cumulative number or effective amount of lineage-specific substitution.

The sensitivity to the ecological barrier is positive. Since
\begin{equation}
\frac{\partial\pi}{\partial m}
=-\frac{2s}{(s+2m)^2}
\end{equation}
and
\begin{equation}
\frac{\partial m_{\mathrm{eff}}}{\partial L_E}
=-m_{\mathrm{eff}},
\end{equation}
we have
\begin{equation}
\boxed{
\frac{\partial\pi}{\partial L_E}
=\frac{2s\,m_{\mathrm{eff}}}{(s+2m_{\mathrm{eff}})^2}>0.}
\label{eq:establishmentBarrierSensitivity}
\end{equation}

\inwords A stronger early barrier reduces effective migration, which increases the fraction of locally favoured substitutions able to establish independently.

\section{Combinatorial DMI accumulation}

If lineage 1 has \(K_1\) derived substitutions and lineage 2 has \(K_2\), there are
\begin{equation}
K_1K_2
\end{equation}
possible cross-lineage pairs. If each pair has aggregate incompatibility probability \(q\), the expected incompatibility count or strength is
\begin{equation}
\boxed{D=qK_1K_2.}
\label{eq:DMIcount}
\end{equation}
For equal lineages \(K_1=K_2=K\),
\begin{equation}
D=qK^2.
\end{equation}

Differentiate the general expression:
\begin{equation}
\dot D=q(\dot K_1K_2+K_1\dot K_2).
\end{equation}
As substitution counts grow, each new substitution has more possible partners in the other lineage. This is the combinatorial origin of Orr-like acceleration.

\section{Map incompatibility count to intrinsic isolation}

Suppose effective incompatibility events act independently and their realized count is Poisson with mean \(\eta D\). The probability of zero effective incompatibilities is
\begin{equation}
\Pr(N=0)=e^{-\eta D}.
\end{equation}
If zero is required for full intrinsic viability, define
\begin{equation}
\boxed{I_D=1-e^{-\eta D}.}
\label{eq:DMIisolation}
\end{equation}

\inwords The exponential mapping is not arbitrary: it is the zero-event probability of a Poisson process. The parameter \(\eta\) converts expected DMI number or strength into fitness effect.

\section{Combine early and late barriers}

If ecological and intrinsic survival components multiply,
\begin{equation}
1-I_{\mathrm{tot}}
=(1-I_E)(1-I_D).
\end{equation}
Using \(1-I_E=e^{-L_E}\) and \(1-I_D=e^{-\eta D}\),
\begin{align}
1-I_{\mathrm{tot}}
&=e^{-L_E}e^{-\eta D}\\
&=e^{-(L_E+\eta D)}.
\end{align}
Therefore
\begin{equation}
\boxed{
I_{\mathrm{tot}}
=1-(1-I_E)(1-I_D)
=1-e^{-[L_E+\eta D]}.}
\label{eq:combinedBarrier}
\end{equation}

\begin{geombox}
Log barriers add because survival probabilities multiply. The early geometric component and the later DMI component occupy different biological layers but share a common additive scale in negative log survival.
\end{geombox}

\section{Early-time expansion and an important nuance}

Suppose the ecological log barrier begins as
\begin{equation}
L_E(t)=bt^2+O(t^3).
\label{eq:earlyEcologicalBarrier}
\end{equation}
Use the scalar exponential expansion \(e^{-x}=1-x+O(x^2)\). Because the unknown remainder in \cref{eq:earlyEcologicalBarrier} is already order \(t^3\),
\begin{align}
m_{\mathrm{eff}}(t)
&=m_0e^{-L_E(t)}\\
&=m_0\left[1-L_E(t)+O(L_E(t)^2)\right]\\
&=m_0\left[1-bt^2+O(t^3)\right].
\label{eq:earlyEffectiveMigrationExpansion}
\end{align}

\inwords The geometric barrier lowers effective migration at order \(t^2\). We retain \(O(t^3)\), rather than claiming \(O(t^4)\), because the ecological barrier itself may contain a cubic term.

It is convenient to regard establishment as a function of the log barrier through \(m_{\mathrm{eff}}=m_0e^{-L_E}\). Taylor-expand around \(L_E=0\):
\begin{align}
\pi_i(L_E)
&=\pi_i(0)
+\left.\frac{\partial\pi_i}{\partial L_E}\right|_0L_E
+O(L_E^2)\\
&=\pi_i(m_0)
+b\left.\frac{\partial\pi_i}{\partial L_E}\right|_0t^2
+O(t^3).
\label{eq:earlyEstablishmentExpansion}
\end{align}
For the closure in \cref{eq:establishmentClosure},
\begin{equation}
\left.\frac{\partial\pi_i}{\partial L_E}\right|_0
=\frac{2s_im_0}{(s_i+2m_0)^2}>0.
\end{equation}

Now integrate \(\dot K_i=\nu_i\pi_i\) from zero to \(t\):
\begin{align}
K_i(t)
&=\int_0^t\nu_i\pi_i(\tau)\,\dd\tau\\
&=\nu_i\pi_i(m_0)t
+\frac{\nu_i b}{3}
\left.\frac{\partial\pi_i}{\partial L_E}\right|_0t^3
+O(t^4).
\label{eq:earlySubstitutionExpansion}
\end{align}
Define
\begin{equation}
a_i=\nu_i\pi_i(m_0),
\qquad
c_i=\frac{\nu_i b}{3}
\left.\frac{\partial\pi_i}{\partial L_E}\right|_0.
\end{equation}
Then
\begin{equation}
K_i(t)=a_it+c_it^3+O(t^4).
\end{equation}
Substitute into \(D=qK_1K_2\):
\begin{align}
D(t)
&=q\left[a_1t+c_1t^3+O(t^4)\right]
\left[a_2t+c_2t^3+O(t^4)\right]\\
&=\boxed{
q a_1a_2t^2
+q(a_1c_2+a_2c_1)t^4
+O(t^5).}
\label{eq:earlyDMIExpansion}
\end{align}

\begin{mathteach}
The leading \(t^2\) term is the ordinary combinatorial product of two substitution counts that each grow linearly from a nonzero baseline establishment rate. The geometric precursor first enters this smooth expansion through \(c_i\), and therefore modifies the \(t^4\) coefficient of \(D\), not its leading exponent.
\end{mathteach}

The intrinsic-isolation mapping adds its own fourth-order correction:
\begin{align}
I_D(t)
&=1-e^{-\eta D(t)}\\
&=\eta D(t)-\half\eta^2D(t)^2+O(D(t)^3).
\end{align}
Thus, even at order \(t^4\), the observed intrinsic-isolation curve combines precursor-driven acceleration with saturation of the exponential mapping.

\begin{biobox}
The robust biological statement is that the early ecological barrier raises the establishment rate of lineage-specific substitutions by reducing homogenizing migration. Under a smooth model with nonzero baseline establishment, this changes the amount of later DMI accumulation before it changes the leading power law. The effect can become much sharper when migration nearly prevents establishment, or when establishment has a threshold, bistability, or other nonlinear feedback.
\end{biobox}

\begin{cautionbox}
The handoff does not automatically create a new leading exponent for DMI accumulation under every smooth establishment function. Its robust claim is causal acceleration through reduced gene flow. Dramatic exponent changes require stronger nonlinearities, thresholds, or feedbacks and should be treated as additional hypotheses.
\end{cautionbox}

\section{Reversibility and the two-stage ratchet}

If environments converge, \(L_E\) can fall as channels and fitness geometries realign. In contrast, substitutions and DMIs that have already accumulated need not disappear. Therefore the early mechanism is potentially reversible, whereas the later component is path dependent and more persistent.

\begin{figure}[H]
\centering
\includegraphics[width=\textwidth]{figures/pedagogy_handoff.pdf}
\caption{Logical handoff from divergent selection to channel rotation, ecological hybrid load, reduced effective gene flow, and later DMI accumulation.}
\label{fig:handoffPedagogy}
\end{figure}

\begin{figure}[H]
\centering
\includegraphics[width=0.84\textwidth]{figures/04_geometric_to_dmi_handoff.pdf}
\caption{Illustrative numerical coupling. The model demonstrates the mechanism rather than fitting empirical data.}
\label{fig:handoffNumerical}
\end{figure}

\part{Connections to earlier quantitative-genetic machinery and empirical tests}

\chapter{The covariance tensor and the new channel geometry}

\section{Two different spaces must not be confused}

The covariance-tensor approach compares whole symmetric matrices such as \(\G_1,\ldots,\G_m\). The channel approach compares invariant subspaces inside ordinary trait space after constructing \(\Sop_i\). These are related but distinct geometries:

\begin{itemize}[leftmargin=1.6em]
\item \textbf{Matrix space:} each complete \(\G_i\) or \(\Sop_i\) is one point.
\item \textbf{Trait space:} each eigenvector of a matrix is a direction; each leading eigenspace is a channel.
\end{itemize}

\begin{cautionbox}
An eigentensor is an axis of variation among matrices. A channel eigenvector is an axis of response among traits. They live in different spaces and answer different questions.
\end{cautionbox}

\section{The space of symmetric matrices}

The real symmetric \(n\times n\) matrices form a vector space \(\operatorname{Sym}(n)\) of dimension
\begin{equation}
q=\frac{n(n+1)}2.
\end{equation}
The Frobenius inner product is
\begin{equation}
\innerF{\mathbf A}{\mathbf B}
=\tr(\mathbf A^\top\mathbf B)
=\tr(\mathbf A\mathbf B)
\end{equation}
for symmetric matrices.

Define \(\svec(\mathbf A)\in\R^q\) by listing diagonal entries without scaling and one copy of every off-diagonal entry multiplied by \(\sqrt2\). For example, for a two-by-two symmetric matrix,
\begin{equation}
\mathbf A=
\begin{bmatrix}
a_{11}&a_{12}\\a_{12}&a_{22}
\end{bmatrix},
\qquad
\svec(\mathbf A)=
\begin{bmatrix}
a_{11}\\a_{22}\\\sqrt2a_{12}
\end{bmatrix}.
\end{equation}

For symmetric \(\mathbf A\) and \(\mathbf B\),
\begin{align}
\svec(\mathbf A)^\top\svec(\mathbf B)
&=\sum_iA_{ii}B_{ii}
+2\sum_{i>j}A_{ij}B_{ij}\\
&=\sum_{i,j}A_{ij}B_{ij}\\
&=\tr(\mathbf A\mathbf B)\\
&=\innerF{\mathbf A}{\mathbf B}.
\label{eq:svecIsometry}
\end{align}

\inwords The \(\sqrt2\) factor makes vectorization preserve matrix lengths and angles. Without it, off-diagonal information would be underweighted because a symmetric off-diagonal entry occurs twice in the full matrix.

\begin{connectionbox}
This is exactly why the existing covariance-tensor code multiplies covariance--covariance blocks by \(2\), cross blocks by \(\sqrt2\), and maps eigentensor off-diagonal loadings back with \(1/\sqrt2\). The code is implementing an isometry between symmetric-matrix space and ordinary Euclidean vector space.
\end{connectionbox}

\section{Construct the covariance tensor as an ordinary covariance matrix}

Let
\begin{equation}
\mathbf g_r=\svec(\G_r),
\qquad r=1,\ldots,m,
\end{equation}
and
\begin{equation}
\bar{\mathbf g}=\frac1m\sum_{r=1}^m\mathbf g_r.
\end{equation}
The covariance of matrix elements is
\begin{equation}
\mathbf C_G
=\frac1{m-1}\sum_{r=1}^m
(\mathbf g_r-\bar{\mathbf g})
(\mathbf g_r-\bar{\mathbf g})^\top.
\label{eq:matrixSpaceCovariance}
\end{equation}
This \(q\times q\) symmetric matrix is the coordinate representation of the fourth-order covariance tensor.

Diagonalize it:
\begin{equation}
\mathbf C_G\mathbf e_s=\alpha_s\mathbf e_s.
\end{equation}
Map each eigenvector back to a symmetric matrix:
\begin{equation}
\mathbf E_s=\svec^{-1}(\mathbf e_s).
\end{equation}
The matrices \(\mathbf E_s\) are eigentensors.

\section{Coordinates of each original matrix}

The coordinate of \(\G_r\) on eigentensor \(\mathbf E_s\) is
\begin{align}
c_{rs}
&=\mathbf e_s^\top\mathbf g_r\\
&=\svec(\mathbf E_s)^\top\svec(\G_r)\\
&=\innerF{\mathbf E_s}{\G_r}\\
&=\tr(\mathbf E_s\G_r).
\label{eq:eigentensorCoordinate}
\end{align}

\inwords A matrix coordinate is a dot product. Large positive and negative coordinates identify matrices that differ strongly and oppositely along an eigentensor axis.

An eigentensor can itself be diagonalized:
\begin{equation}
\mathbf E_s=\sum_{j=1}^n\eta_{sj}\mathbf e_{sj}\mathbf e_{sj}^\top.
\end{equation}
The vectors \(\mathbf e_{sj}\) identify trait combinations that contribute to the matrix-space difference captured by \(\mathbf E_s\).

\begin{geombox}
The first eigenanalysis finds directions through the cloud of whole matrices. The second eigenanalysis opens one such direction back up and asks which trait combinations create it.
\end{geombox}

\section{Three uses in the geometric snowball programme}

\subsection{Tensor on \(\G_i\)}
This reproduces the earlier question: how has the distribution of additive variance changed among populations or ecotypes? It is directly continuous with the established \textit{Senecio} analyses \cite{walter2018,aguirre2014,hine2009}.

\subsection{Tensor on \(\Hh_i\)}
This asks how local fitness curvature differs among habitats. It can identify recurrent combinations of traits in which selection surfaces diverge.

\subsection{Tensor on \(\Sop_i\)}
This asks how the complete joint response operators differ. It is useful for describing operator-level divergence across many populations. Principal-angle analysis is then applied to the leading eigenspaces of each \(\Sop_i\) to quantify channel mismatch directly.

\begin{figure}[H]
\centering
\includegraphics[width=\textwidth]{figures/pedagogy_tensor_bridge.pdf}
\caption{Bridge from the established covariance-tensor framework to response-operator and channel geometry. Eigentensors describe axes of whole-matrix divergence; principal angles describe divergence among response subspaces.}
\label{fig:tensorBridge}
\end{figure}

\section{A new synthesis statistic}

Across multiple populations, one can separate three levels:
\begin{enumerate}[leftmargin=1.6em]
\item \(\alpha_s\): how much matrices differ along eigentensor \(s\);
\item \(c_{rs}\): which populations occupy opposite ends of that matrix-space axis;
\item \(\psi_{ij,1},\ldots,\psi_{ij,k}\): the principal angles measuring how much the biologically chosen response subspaces of populations \(i\) and \(j\) differ.
\end{enumerate}

This makes it possible to ask whether the major axes of operator divergence identified by a tensor are the same changes that generate large channel angles and hybrid load.

\chapter{Mean-divergence matrices and their relation to channels}

\section{The \(\mathbf D\) matrix}

Earlier \textit{Senecio} work summarized divergence in multivariate mean phenotype using a covariance-like matrix \(\mathbf D\). In a balanced nested design, a form such as
\begin{equation}
\mathbf D=\frac{\mathbf{MS}_{\mathrm{ecotype}}-\mathbf{MS}_{\mathrm{population(ecotype)}}}{n_f}
\end{equation}
separates between-ecotype mean divergence from population variation within ecotypes \cite{walter2018}.

Its eigendecomposition
\begin{equation}
\mathbf D=\sum_jd_j\mathbf d_j\mathbf d_j^\top
\end{equation}
identifies directions of realized mean divergence.

\section{Projection of mean divergence onto a channel}

For a unit direction \(\mathbf u\), the amount of mean divergence associated with that direction is
\begin{equation}
V_D(\mathbf u)=\mathbf u^\top\mathbf D\mathbf u.
\label{eq:Dprojection}
\end{equation}
If \(\mathbf u=\sum_jc_j\mathbf d_j\), then
\begin{align}
\mathbf u^\top\mathbf D\mathbf u
&=\sum_jd_jc_j^2.
\end{align}

\inwords Projection asks how much of the realized mean divergence lies along a proposed trait combination. It is a directional-variance calculation analogous to \(\mathbf u^\top\G\mathbf u\).

\section{What \(\mathbf D\) can and cannot tell us}

\(\mathbf D\) can test whether mean divergence aligns with ancestral or derived channels. It does not estimate fitness curvature and cannot replace \(\Hh\). Realized mean divergence reflects optima, historical forcing, gene flow, and response geometry together.

\begin{cautionbox}
Phenotypic divergence \(\Delta\muv\), environmental Euclidean distance \(\Delta E\), and curvature difference \(\Delta\Hh\) are not interchangeable. Environmental variables describe habitats; \(\Hh\) describes how fitness changes with phenotype inside a habitat. A large environmental distance can produce little curvature difference in the measured traits, and a small environmental distance can rotate selection strongly.
\end{cautionbox}

\section{A sequence of increasingly strong tests}

\begin{enumerate}[leftmargin=1.6em]
\item Does \(\mathbf D\) align with \(\mathbf g_{\max}\)?
\item Does \(\mathbf D\) align better with the leading channel of \(\Sop\)?
\item Do between-population channel angles predict recombinant hybrid load?
\item Does the exact curvature-weighted Gaussian model outperform raw channel angle?
\end{enumerate}

Each step retains the earlier result as a nested or limiting model.

\chapter{What the existing \textit{Senecio} results already contribute}

\section{Local adaptation and early reproductive isolation}

Reciprocal transplants across Dune, Headland, Tableland, and Woodland environments showed strong local adaptation and substantial immigrant inviability. F1 hybrids exhibited heterosis, yet hybrids with a native parent performed better than hybrids with two foreign parents. F2 hybrids showed a pronounced fitness reduction in a common garden, consistent with early intrinsic incompatibilities emerging after recombination \cite{walter2016}.

\begin{biobox}
This combination is exactly the biological setting in which an early geometric mechanism is plausible: strong environment-dependent barriers are already present, F1 vigor can mask incompatibility, and recombinant generations reveal additional failure.
\end{biobox}

\section{Evolution of genetic variance}

The covariance-tensor analysis of ecotype \(\G\) matrices found divergence in the magnitude and multivariate distribution of genetic variance. Divergence in \(\G\), including divergence in the orientation of \(\mathbf g_{\max}\), aligned with divergence in mean phenotype \cite{walter2018}.

\begin{connectionbox}
That result supplies empirical precedent for the first half of the geometric snowball: divergent natural selection can be associated with changes in the distribution and orientation of standing genetic variance. The new theory adds local fitness curvature and asks whether the resulting \(\Sop\) subspaces predict hybrid load.
\end{connectionbox}

\section{Loss of ecologically important variation after hybrid breakdown}

In later-generation hybrids, severe F2 fertility loss was followed by fitness recovery, reduced variance in field performance, loss of parental fitness trade-offs, and loss of phenotypic extremes in the F4. These patterns suggest that variation associated with adaptation was also involved in negative interactions contributing to intrinsic isolation \cite{walter2020}.

\begin{connectionbox}
The trace identity gives a continuous quantitative version of this idea. It asks whether covariance adapted to one curvature matrix becomes costly when averaged or recombined with covariance adapted to another. The DMI layer then captures additional locus-specific negative interactions that need not be visible in additive covariance alone.
\end{connectionbox}

\section{The replicated Dune--Headland opportunity}

The strongest test is pair-specific rather than ecotype-pooled. Replicate Dune--Headland pairs allow environmental repeatability to be separated from pair-specific mutation order. The theory predicts that repeatable habitat forcing should contribute a repeatable component to \(\dot\Sop\) and \(\cK\), whereas stochastic histories should generate pair-specific departures.

\begin{figure}[H]
\centering
\includegraphics[width=0.90\textwidth]{figures/senecio_pair_descriptive.pdf}
\caption{Exploratory pair-level summary from available morphology and soil data. Both axes are descriptive distances only. The weak association is informative because it reinforces that soil-chemistry distance cannot substitute for directly estimated fitness curvature and morphology-mean distance cannot substitute for channel angle.}
\label{fig:senecioExploratory}
\end{figure}

\chapter{Estimating the theory from data}

\section{Estimate \(\G_i\) at the population-pair level}

A multivariate animal model or breeding design can estimate posterior draws
\begin{equation}
\G_i^{(s)},
\qquad s=1,\ldots,S.
\end{equation}
Population-pair-specific matrices are preferable because ecotype pooling can average over precisely the among-pair channel variation the theory seeks to explain.

\section{Estimate local curvature from phenotype--fitness data}

Near an optimum, fit a regularized quadratic selection model
\begin{equation}
\log W
=\alpha+\betav^\top\mathbf z
-\half\mathbf z^\top\Hh\mathbf z+\varepsilon.
\label{eq:empiricalQuadraticSelection}
\end{equation}
If traits are centred at the estimated optimum, the linear term should be small. If they are not centred, the optimum implied by the quadratic form is
\begin{equation}
\thetav=\Hh^{-1}\betav,
\end{equation}
provided \(\Hh\succ0\). This follows because
\begin{align}
\betav^\top\mathbf z-\half\mathbf z^\top\Hh\mathbf z
&=-\half(\mathbf z-\Hh^{-1}\betav)^\top\Hh
(\mathbf z-\Hh^{-1}\betav)
+\half\betav^\top\Hh^{-1}\betav.
\end{align}

\begin{cautionbox}
Quadratic regression is parameter hungry: \(n(n+1)/2\) curvature parameters are required. Shrinkage, structured low-rank models, preregistered trait modules, or Bayesian regularization are likely necessary. Positive-definiteness should be imposed or carefully assessed rather than assumed from noisy estimates.
\end{cautionbox}

\section{Construct posterior channel quantities}

For each posterior draw:
\begin{enumerate}[leftmargin=1.6em]
\item compute \(\G_i^{1/2}\);
\item compute \(\Sop_i=\G_i^{1/2}\Hh_i\G_i^{1/2}\);
\item choose or average over channel dimension \(k\);
\item map channel bases to phenotype space;
\item construct common-metric bases \(\Qq_i\);
\item compute principal angles, \(d_C^2\), \(I_V\), and spectral gaps;
\item compute exact expected log load and Gaussian arithmetic mean fitness for hybrids;
\item retain the full posterior joint distribution rather than plugging in posterior means.
\end{enumerate}

\section{Estimate hybrid moments}

For each cross and environment, estimate
\begin{equation}
\muv_H,
\qquad
\Pp_H.
\end{equation}
Then evaluate
\begin{equation}
L^{\log}_{H,e}
=\half(\muv_H-\thetav_e)^\top\Hh_e(\muv_H-\thetav_e)
+\half\tr(\Hh_e\Pp_H)
\end{equation}
and, where Gaussian approximation is defensible, \cref{eq:exactGaussianMeanFitness}.

\section{Competing predictive models}

A preregistered comparison should include:
\begin{enumerate}[leftmargin=1.6em]
\item divergence time or neutral genetic distance alone;
\item environmental distance alone;
\item mean mismatch alone;
\item \(\G\)-matrix distance alone;
\item \(\Hh\)-matrix distance alone;
\item channel angle or chordal distance alone;
\item the full curvature-weighted hybrid-load model;
\item the coupled ecological-plus-DMI model for later generations.
\end{enumerate}

The key evidentiary standard is out-of-sample prediction of hybrid fitness, not merely a significant association in the fitted data.

\chapter{Computational algorithms}

\section{Algorithm 1: one-population channel decomposition}

\begin{lstlisting}
Input: SPD matrices G and H; channel dimension k.
1. Symmetrise numerical estimates: G <- (G + G^T)/2, H <- (H + H^T)/2.
2. Eigendecompose G = Q diag(g) Q^T.
3. Check all g_j > tolerance.
4. Form Ghalf = Q diag(sqrt(g)) Q^T.
5. Form S = Ghalf H Ghalf and symmetrise.
6. Eigendecompose S; order eigenvalues from largest to smallest.
7. Retain V = first k eigenvectors.
8. Map to phenotype space U = Ghalf V.
Output: eigenvalues, whitened basis V, phenotype basis U, spectral gap delta_k.
\end{lstlisting}

\section{Algorithm 2: invariant pairwise channel comparison}

\begin{lstlisting}
Input: phenotype bases U1 and U2; common SPD metric M.
1. A1 = U1^T M U1; A2 = U2^T M U2.
2. Q1 = U1 A1^{-1/2}; Q2 = U2 A2^{-1/2}.
3. Compute singular values sigma of Q1^T M Q2.
4. Clip numerical sigma to [0,1].
5. psi_j = arccos(sigma_j).
6. dC2 = sum sin(psi_j)^2.
7. IV = 1 - product cos(psi_j)^2.
Output: principal angles, chordal distance, volume deficit.
\end{lstlisting}

\section{Algorithm 3: operator-eigenspace susceptibility in a fixed representation}

\begin{lstlisting}
Input: ancestral S0; operator velocities A1 and A2; channel dimension k.
1. Eigendecompose S0 = [V W] diag(lambda) [V W]^T.
2. For every included b and excluded a:
       KS[a,b] = W[:,a]^T (A2-A1) V[:,b] /
                 (lambda[b]-lambda[k+a]).
3. chiS2 = sum KS[a,b]^2.
Output: KS and chiS2. This is the full coefficient when the whitening
        map is shared and fixed; otherwise continue with Algorithm 3b.
\end{lstlisting}

\section{Algorithm 3b: fully invariant susceptibility when \(\G\) evolves}

\begin{lstlisting}
Input: G0, H0, dG1, dH1, dG2, dH2, common metric M, dimension k.
1. R0 = sqrt(G0); S0 = R0 H0 R0.
2. Eigendecompose S0 and retain V, W, and eigenvalues lambda.
3. Solve R0 Xi + Xi R0 = dGi for X1 and X2.
4. For each population form
       dSi = Xi H0 R0 + R0 dHi R0 + R0 H0 Xi.
5. Compute KSi from dSi using the eigengap formula in Algorithm 3.
6. U0 = R0 V; B0 = sqrt(M) U0; Gamma0 = B0^T B0.
7. Q0 = B0 Gamma0^{-1/2}; construct an orthonormal complement Qperp.
8. KM = Qperp^T sqrt(M) [(X2-X1)V + R0 W(KS2-KS1)]
        Gamma0^{-1/2}.
9. chiM2 = sum KM[a,b]^2.
Output: KM and chiM2, where dC,M^2(t) = chiM2 t^2 + O(t^3).
\end{lstlisting}

\section{Algorithm 4: exact hybrid fitness}

\begin{lstlisting}
Input: hybrid mean mu, covariance P, environment optimum theta, curvature H.
1. d = mu - theta.
2. Expected log load:
       Llog = 0.5 d^T H d + 0.5 trace(H P).
3. For Gaussian arithmetic mean fitness:
       Hhalf = sqrt(H)
       C = Hhalf P Hhalf
       m = Hhalf d
       Wrel = det(I+C)^(-1/2) *
              exp[-0.5 m^T (I+C)^(-1) m].
Output: Llog and Wrel.
\end{lstlisting}

\section{Numerical stability}

\begin{itemize}[leftmargin=1.6em]
\item Use symmetric eigensolvers, not general eigensolvers, for \(\G\), \(\Hh\), and \(\Sop\).
\item Symmetrise estimates before decomposition to remove floating-point asymmetry.
\item Do not invert matrices explicitly when a linear solve is available.
\item Treat small or negative posterior eigenvalues carefully; they can indicate estimation noise or an invalid curvature model.
\item Compare subspaces rather than individual vectors near eigenvalue ties.
\item Propagate posterior uncertainty through every nonlinear operation.
\end{itemize}

\section{Independent numerical audit}

The companion script checks whitening, similarity of \(\G\Hh\) and \(\Sop\), generalized eigenvector residuals, common-metric coordinate invariance, projector and determinant identities, the fixed-representation perturbation coefficient, commuting perturbations, the full common-metric coefficient for evolving \(\G\) and \(\Hh\), coordinate invariance of that coefficient, both hybrid-fitness formulas, the general trace identity, the ridge theorem, the matrix-square-root derivative, the standing-variation response with migration, and the combined-barrier identity.

All nineteen checks pass under the fixed audit seed. The complete log is reproduced in \cref{app:numericalChecks}.

\chapter{Falsification criteria and boundary cases}

\section{Results that would weaken or reject the central mechanism}

The geometric mechanism would be weakened if:
\begin{enumerate}[leftmargin=1.6em]
\item directly estimated \(\Sop\) channels do not diverge more between contrasting habitats than within habitats;
\item channel divergence is not repeatably associated with divergent selection;
\item curvature-weighted hybrid covariance fails to predict recombinant hybrid fitness out of sample;
\item raw neutral divergence performs as well as or better than the full geometric model;
\item hybrid load persists identically across environments despite the model predicting a large extrinsic component;
\item estimated covariance--curvature alignment has the wrong sign but hybrids nevertheless show the predicted positive load without another mechanism;
\item inferred channel angle is driven entirely by an unstable rank-one choice and disappears for the supported channel dimension.
\end{enumerate}

\section{Biologically meaningful cases in which the mechanism predicts little isolation}

\subsection{Commuting divergence}
\(\G\) and \(\Hh\) change in magnitude but retain common axes. Eigenvalues change; channels do not rotate.

\subsection{Large gap and weak cross-forcing}
The channel is strongly protected and the perturbation barely couples it to other modes.

\subsection{Rotated channels without covariance--curvature coadaptation}
Geometric discrepancy exists, but the trace term is zero or negative.

\subsection{Hybrid variance remains on benign directions}
Segregation does not populate directions strongly penalised by either habitat.

\subsection{Dominance masks early additive load}
F1 hybrids are vigorous even though F2 recombination later reveals a cost.

\section{Scope of the local theory}

The perturbation expansion assumes smooth operator paths and a nonzero gap at the chosen boundary. Near an exact crossing, the correct object is the combined degenerate subspace; rank-one derivatives diverge because the rank-one direction is not identifiable. At large divergence, global geometry, multiple optima, nonquadratic selection, developmental constraints, genotype--environment interactions, and epistasis can all matter.

\begin{summarybox}
The local theory is strongest as an early-speciation model. That is a feature rather than a weakness: it targets the interval in which ecological barriers can arise before many fixed intrinsic incompatibilities exist.
\end{summarybox}

\chapter{Integrated predictions}

\begin{longtable}{@{}p{0.06\textwidth}p{0.25\textwidth}p{0.30\textwidth}p{0.30\textwidth}@{}}
\caption{Predictions, mathematical sources, and biological tests.}\\
\toprule
ID & Prediction & Mathematical source & Empirical test \\
\midrule
\endfirsthead
\toprule
ID & Prediction & Mathematical source & Empirical test \\
\midrule
\endhead
P1 & Divergence follows leading modes of \(\Sop\), not necessarily \(\mathbf g_{\max}\). & \(\Sop=\G^{1/2}\Hh\G^{1/2}\); Schluter limit when \(\Hh\propto\Ii\). & Compare realized response with \(\mathbf g_{\max}\) and \(\Sop\) channels. \\
P2 & Deterministic channel mismatch is quadratic at early time. & \(d_{C,M}^2=t^2\normF{\mathbf K_M}^2+O(t^3)\); \(\mathbf K_M\) includes both eigenspace and changing-\(\G^{1/2}\) terms. & Repeated selection trajectories or replicated young population pairs with posterior \(\G\) and \(\Hh\) dynamics. \\
P3 & Drift contributes a linear early baseline. & \(\E[d_C^2]=a t+O(t^2)\); \(a=\tr(\boldsymbol\Sigma_K\boldsymbol\Sigma_K^\top)\). & Estimate the linear term and directional repeatability across replicates; model second-order diffusion before attributing the quadratic coefficient entirely to selection. \\
P4 & Spectral gap alone is insufficient. & \(\cK_{ab}=\mathbf w_a^\top\Delta\Aa\mathbf v_b/(\lambda_b-\lambda_{k+a})\). & Test cross-channel forcing and gap jointly. \\
P5 & Channel angle predicts fitness only under calibration. & General trace identity and ridge theorem. & Compare raw angle against curvature-weighted hybrid covariance. \\
P6 & Recombinant generations reveal the covariance mechanism more clearly than F1. & \(\Rr\) and dominance/segregation distinction. & Assay F1, F2, and advanced hybrids in common and parental environments. \\
P7 & Early isolation is environment dependent and potentially reversible. & Dependence of \(L_E\) on \(\Hh_e\). & Relax divergent selection or increase gene flow experimentally. \\
P8 & Strong early barriers facilitate later intrinsic isolation. & \(m_{\mathrm{eff}}=m_0e^{-L_E}\), \(D=qK_1K_2\). & Conditional on age, test whether early load predicts lower gene flow and stronger common-garden breakdown. \\
P9 & Matrix-space divergence and subspace divergence are related but nonidentical. & Tensor on \(\G\), \(\Hh\), or \(\Sop\) versus principal angles. & Joint tensor and Grassmannian analysis across many populations. \\
\bottomrule
\end{longtable}

\chapter{The complete causal statement}

The most defensible formulation is:

\begin{quote}
Divergent natural selection can change both the curvature of local fitness surfaces and the distribution of standing genetic variance. These changes rotate the leading invariant subspaces of the symmetric response operator \(\Sop=\G^{1/2}\Hh\G^{1/2}\) and, when \(\G\) evolves, also deform the map from whitened coordinates back to phenotype space. Under smooth deterministic divergence, common-metric squared channel distance increases quadratically at early time, with an exact coefficient \(\normF{\mathbf K_M}^2\) determined by cross-channel forcing, spectral gaps, and the changing square-root maps. Channel distance becomes reproductive isolation only when parental covariance structures are coadapted to their own fitness curvatures, so that recombinant hybrids express variance in directions strongly penalised in parental habitats. This early ecological barrier can reduce effective gene flow and thereby facilitate the later combinatorial accumulation of intrinsic Dobzhansky--Muller incompatibilities.
\end{quote}

The theory is therefore neither a replacement for Lande's response equation nor a replacement for Orr's snowball. It is a mathematical bridge between them.

\begin{figure}[H]
\centering
\includegraphics[width=0.96\textwidth]{figures/figure_geometric_snowball_concept.pdf}
\caption{Integrated conceptual view. Ecological divergence rotates parental channels; the resulting overlap deficit begins slowly and then accelerates. The fitness interpretation requires the covariance--curvature calibration derived in this document.}
\label{fig:integratedConcept}
\end{figure}


\appendix

\chapter{Notation glossary}

\begin{longtable}{@{}p{0.18\textwidth}p{0.72\textwidth}@{}}
\toprule
Symbol & Meaning \\
\midrule
\endfirsthead
\toprule
Symbol & Meaning \\
\midrule
\endhead
\(\mathbf z,\zbar\) & Individual phenotype vector and population mean phenotype. \\
\(\thetav\) & Local phenotypic optimum. \\
\(\mathbf d=\zbar-\thetav\) & Mean displacement from the optimum. \\
\(\betav\) & Directional selection gradient, locally \(-\Hh\mathbf d\). \\
\(\G\) & Additive genetic variance--covariance matrix. \\
\(\Hh\) & Curvature matrix of minus log fitness near a local maximum. \\
\(\Sop=\G^{1/2}\Hh\G^{1/2}\) & Symmetric channel or response operator. \\
\(\mathbf v_j\) & Orthonormal eigenvector of \(\Sop\) in whitened coordinates. \\
\(\mathbf u_j=\G^{1/2}\mathbf v_j\) & Corresponding phenotype-space response direction. \\
\(\Vv\) & Matrix whose columns form a channel basis. \\
\(\Pp=\Vv\Vv^\top\) & Orthogonal projector onto a channel subspace. \\
\(k\) & Chosen channel dimension. \\
\(\delta_k=\lambda_k-\lambda_{k+1}\) & Spectral gap separating the channel from its complement. \\
\(\psi_j\) & Principal angles between two rank-\(k\) channel subspaces. \\
\(d_C^2\) & Squared chordal distance, \(\sum_j\sin^2\psi_j\). \\
\(I_V\) & Bounded volume-overlap deficit, \(1-\prod_j\cos^2\psi_j\). \\
\(\Aa_i=\dot\Sop_i(0)\) & Initial operator velocity for population \(i\). \\
\(\cK\) & Relative cross-channel tangent forcing matrix in a fixed shared representation. \\
\(\Delta\mathbf K_S\) & Relative rotation tangent of the leading eigenspace of \(\Sop\). \\
\(\mathbf X_i=D(\G^{1/2})[\dot\G_i]\) & Derivative of the square-root map, obtained from a Sylvester equation. \\
\(\mathbf K_M\) & Full common-metric phenotype-channel tangent, including both \(\Delta\mathbf K_S\) and direct deformation by changing \(\G^{1/2}\). \\
\(\normF{\mathbf K_M}^2\) & Exact invariant local deterministic susceptibility when \(\G\) and \(\Hh\) evolve. \\
\(\M\) & Common SPD metric used to compare phenotype-space channels. \\
\(\Pp_H\) & Hybrid phenotypic covariance. \\
\(\Rr\) & Hybrid covariance excess beyond the parental average; also used contextually for \(\G^{1/2}\) in the square-root chapter, where its meaning is explicitly stated. \\
\(L^{\log}_{H,e}\) & Expected negative log relative fitness of hybrids in environment \(e\). \\
\(L_{\mathrm{cov}}\) & Hybrid covariance log-load relative to parental benchmarks. \\
\(L_E\) & Total early ecological/geometric log barrier. \\
\(m_{\mathrm{eff}}\) & Effective migration after early selection. \\
\(K_i\) & Lineage-specific substitution count or effective amount. \\
\(D=qK_1K_2\) & Expected number or aggregate strength of cross-lineage DMIs. \\
\(I_D\) & Intrinsic DMI isolation. \\
\bottomrule
\end{longtable}

\begin{cautionbox}[title={One notation collision deliberately avoided in computation}]
The literature often uses \(\mathbf R\) both for a matrix square root and for segregation excess. In the prose above, context distinguishes them, but computational implementations should use unambiguous names such as \texttt{Ghalf} and \texttt{Rseg}.
\end{cautionbox}

\chapter{A numerical two-trait example from beginning to end}

\section{Specify \(\G\) and two curvature matrices}

Let
\begin{equation}
\G=
\begin{bmatrix}
2&0\\0&0.5
\end{bmatrix},
\end{equation}
so
\begin{equation}
\G^{1/2}=
\begin{bmatrix}
\sqrt2&0\\0&1/\sqrt2
\end{bmatrix}.
\end{equation}
Consider two populations with
\begin{equation}
\Hh_1=
\begin{bmatrix}
1&0.3\\0.3&1.4
\end{bmatrix},
\qquad
\Hh_2=
\begin{bmatrix}
1&-0.3\\-0.3&1.4
\end{bmatrix}.
\end{equation}
Both are positive definite because their determinants are
\begin{equation}
1(1.4)-0.3^2=1.31>0
\end{equation}
and their first diagonal entry is positive.

\section{Construct the channel operators}

Using \cref{eq:S2D},
\begin{equation}
\Sop_1=
\begin{bmatrix}
2&0.3\\0.3&0.7
\end{bmatrix},
\qquad
\Sop_2=
\begin{bmatrix}
2&-0.3\\-0.3&0.7
\end{bmatrix}.
\end{equation}
The eigenvalues of either matrix are
\begin{align}
\lambda_{1,2}
&=\frac{2+0.7}{2}
\pm\half\sqrt{(2-0.7)^2+4(0.3)^2}\\
&=1.35\pm\half\sqrt{2.05}\\
&\approx2.0659,\;0.6341.
\end{align}
The gap is
\begin{equation}
\delta\approx1.4318.
\end{equation}

For population 1,
\begin{equation}
\phi_1
=\half\operatorname{atan2}(0.6,1.3)
\approx12.39^\circ.
\end{equation}
For population 2,
\begin{equation}
\phi_2\approx-12.39^\circ.
\end{equation}
Thus the whitened channel angle is
\begin{equation}
\theta_{12}=|\phi_1-\phi_2|
\approx24.78^\circ.
\end{equation}

\section{Compute geometric discrepancy}

For \(k=1\),
\begin{equation}
d_C^2=I_V=\sin^2\theta_{12}.
\end{equation}
Numerically,
\begin{equation}
d_C^2\approx\sin^2(24.78^\circ)\approx0.176.
\end{equation}
The overlap is therefore about \(1-0.176=0.824\) on this squared projection scale.

\section{Map to phenotype-space directions}

From \cref{eq:phenotypeAngle2D},
\begin{equation}
\tan\alpha_i
=\sqrt{\frac{0.5}{2}}\tan\phi_i
=0.5\tan\phi_i.
\end{equation}
Therefore
\begin{equation}
\alpha_1\approx6.26^\circ,
\qquad
\alpha_2\approx-6.26^\circ.
\end{equation}
The ordinary phenotype-space angle is smaller because mapping back stretches the high-variance first axis more strongly.

\section{Calibrate a hypothetical ridge load}

Suppose an associated ridge model has
\begin{equation}
\Delta g=1.5,
\qquad
\Delta h=1.0.
\end{equation}
Then
\begin{align}
L_{\mathrm{cov}}
&=\frac{1.5\times1.0}{4}(0.176)\\
&\approx0.066.
\end{align}
The corresponding bounded log-scale isolation is
\begin{equation}
I_L=1-e^{-0.066}\approx0.064.
\end{equation}

\begin{cautionbox}
The number \(0.176\) is pure geometry. The number \(0.066\) is a fitness load only because \(\Delta g\) and \(\Delta h\) have supplied a biological scale. Different covariance and curvature contrasts would map the same angle to a different load.
\end{cautionbox}

\chapter{Useful matrix identities collected in one place}

\begin{longtable}{@{}p{0.31\textwidth}p{0.59\textwidth}@{}}
\toprule
Identity & Interpretation and use \\
\midrule
\endfirsthead
\toprule
Identity & Interpretation and use \\
\midrule
\endhead
\(\tr(\mathbf A\mathbf B)=\tr(\mathbf B\mathbf A)\) & Cyclic trace invariance; moves factors to expose covariance or projector products. \\
\(\E[\boldsymbol\varepsilon^\top\mathbf A\boldsymbol\varepsilon]=\tr(\mathbf A\Pp)\) & Converts quadratic variance into expected loss. \\
\(\normF{\mathbf A}^2=\tr(\mathbf A^\top\mathbf A)\) & Matrix analogue of squared Euclidean length. \\
\(\half\normF{\Pp_1-\Pp_2}^2=\sum_j\sin^2\psi_j\) & Chordal distance between equal-rank subspaces. \\
\(\tr(\Pp_1\Pp_2)=\sum_j\cos^2\psi_j\) & Total squared projection overlap. \\
\(\det(\Vv_1^\top\Vv_2)^2=\prod_j\cos^2\psi_j\) & Squared projected \(k\)-volume overlap. \\
\(\G\Hh=\G^{1/2}\Sop\G^{-1/2}\) & Similarity between original and symmetric response operators. \\
\(\Hh\mathbf u=\lambda\G^{-1}\mathbf u\) & Generalized symmetric eigenproblem for phenotype channels. \\
\(\Rr\mathbf X+\mathbf X\Rr=\dot\G\) & Fr\'echet derivative of the matrix square root. \\
\(d_{C,M}^2(t)=t^2\normF{\mathbf K_M}^2+O(t^3)\) & Fully invariant local snowball when the phenotype-space channel and \(\G^{1/2}\) both move. \\
\(\svec(\mathbf A)^\top\svec(\mathbf B)=\tr(\mathbf A\mathbf B)\) & Isometric vectorization used by covariance tensors. \\
\(1-(1-I_1)(1-I_2)=1-e^{-(L_1+L_2)}\) & Independent survival barriers multiply; log barriers add. \\
\bottomrule
\end{longtable}

\chapter{Numerical audit output}
\label{app:numericalChecks}

\lstinputlisting{checks/derivation_audit.txt}

\chapter{What was added or corrected relative to the shorter theory note}

\begin{enumerate}[leftmargin=1.6em]
\item The local Taylor expansion, selection gradient, whitening transformation, symmetric operator, generalized eigenproblem, and mode solution are derived line by line.
\item The fixed-\(\G\) assumption behind \(\dot{\mathbf x}=-\Sop\mathbf x\) is stated explicitly; the moving-metric term is identified when \(\G(t)\) changes.
\item Channel comparison for \(\G_1\neq\G_2\) is proven coordinate invariant under a common metric, while the scientific choice of that metric is distinguished from mere coordinate invariance.
\item The projector proof of the fixed-representation local quadratic law is supplied in full.
\item A second theorem derives the fully invariant coefficient \(\mathbf K_M\) when \(\G\) evolves, including both rotation of the \(\Sop\) eigenspace and direct deformation by \(\G^{1/2}\).
\item The role of commutation and off-diagonal forcing is separated from the role of the spectral gap; a diagonal example shows explicitly that a positive Davis--Kahan bound can accompany zero actual rotation.
\item The deterministic-versus-stochastic result is separated into an exact tangent-surrogate identity \(a t+b t^2\) and the actual chordal asymptotic \(a t+O(t^2)\), whose second-order term can contain stochastic curvature corrections.
\item Expected log fitness and arithmetic mean fitness are distinguished and each is derived separately.
\item The general trace identity and channel-ridge theorem are proven in full, including the condition required for the ridge to be the leading response channel.
\item The standing-variation derivation now includes the migration contribution
\[
\dot{\zbar}_i=\G_i\betav_i+m(\zbar_j-\zbar_i).
\]
\item The derivative of \(\G^{1/2}\) is solved explicitly in the eigenbasis of \(\G\), allowing \(\dot\Sop\) to be partitioned into genetic and curvature components.
\item The DMI handoff is given a Poisson interpretation and its early-time expansion clarifies what it does and does not change about the leading exponent.
\item The existing covariance-tensor machinery is derived through the isometric \(\svec\) map and connected explicitly to the new response-operator geometry.
\end{enumerate}

\clearpage
\begin{thebibliography}{99}

\bibitem{aguirre2014}
Aguirre, J. D., Hine, E., McGuigan, K. and Blows, M. W. (2014).
Comparing \(\G\): multivariate analysis of genetic variation in multiple populations.
\emph{Heredity} \textbf{112}, 21--29.

\bibitem{bjorckgolub1973}
Bj\"orck, \AA. and Golub, G. H. (1973).
Numerical methods for computing angles between linear subspaces.
\emph{Mathematics of Computation} \textbf{27}, 579--594.

\bibitem{daviskahan1970}
Davis, C. and Kahan, W. M. (1970).
The rotation of eigenvectors by a perturbation. III.
\emph{SIAM Journal on Numerical Analysis} \textbf{7}, 1--46.

\bibitem{hine2009}
Hine, E., Chenoweth, S. F., Rundle, H. D. and Blows, M. W. (2009).
Characterizing the evolution of genetic variance using genetic covariance tensors.
\emph{Philosophical Transactions of the Royal Society B} \textbf{364}, 1567--1578.

\bibitem{lande1979}
Lande, R. (1979).
Quantitative genetic analysis of multivariate evolution, applied to brain:body size allometry.
\emph{Evolution} \textbf{33}, 402--416.

\bibitem{landearnold1983}
Lande, R. and Arnold, S. J. (1983).
The measurement of selection on correlated characters.
\emph{Evolution} \textbf{37}, 1210--1226.

\bibitem{lynchwalsh1998}
Lynch, M. and Walsh, B. (1998).
\emph{Genetics and Analysis of Quantitative Traits}.
Sinauer Associates, Sunderland, Massachusetts.

\bibitem{orr1995}
Orr, H. A. (1995).
The population genetics of speciation: the evolution of hybrid incompatibilities.
\emph{Genetics} \textbf{139}, 1805--1813.

\bibitem{schluter1996}
Schluter, D. (1996).
Adaptive radiation along genetic lines of least resistance.
\emph{Evolution} \textbf{50}, 1766--1774.

\bibitem{walshlynch2018}
Walsh, B. and Lynch, M. (2018).
\emph{Evolution and Selection of Quantitative Traits}.
Oxford University Press, Oxford.

\bibitem{walter2016}
Walter, G. M., Wilkinson, M. J., James, M. E., Richards, T. J., Aguirre, J. D. and Ortiz-Barrientos, D. (2016).
Diversification across a heterogeneous landscape.
\emph{Evolution} \textbf{70}, 1979--1992.

\bibitem{walter2018}
Walter, G. M., Aguirre, J. D., Blows, M. W. and Ortiz-Barrientos, D. (2018).
Evolution of genetic variance during adaptive radiation.
\emph{The American Naturalist} \textbf{191}, E108--E128.

\bibitem{walter2020}
Walter, G. M., Richards, T. J., Wilkinson, M. J., Blows, M. W., Aguirre, J. D. and Ortiz-Barrientos, D. (2020).
Loss of ecologically important genetic variation in late generation hybrids reveals links between adaptation and speciation.
\emph{Evolution Letters} \textbf{4}, 302--316.

\end{thebibliography}


\end{document}
